
\chapter{导出函子}
\label{ch:derived-functors}

\epigraph{把这么一大堆东西——单独拿出来都算不上有趣的东西——统统塞进「导出函子」这项大帽子底下，在我看来实在是件乐事。}{——\trGrothendieck，致 \trSerre 书信，1955 年 2 月 18 日 [GrS, p.~6]}

本章讨论\trDerivedFunctor——由\trGrothendieck 在其著名的「东北论文」[Gr1] 中引入——以及它们在\trSheaf{}上的应用。对拟紧分离概形上的\trQCSheaf，导出函子上同调将与\trCechCohomology{}一致（\S23.5）。\trCechCohomology{}对我们的大多数目的而言已经足够，而且相当接地气、可以动手计算；但导出函子上同调仍然值得一见。首先，它在代数几何中的应用要广泛得多（比如 \trEtale{}上同调），在别的领域也是如此，尽管这超出了本书的范围。其次，它将轻而易举地为我们提供一些有用的概念，比如 $\mathrm{Ext}$ 函子与\trLeray{}谱序列。不过，导出函子头一回见面可能令人生畏，所以你尽管先把主要结果浏览一遍，日后再回头细读。


\section{$\mathrm{Tor}$ 函子}

我们从一道热身开始：$\mathrm{Tor}$ 的情形。这是一个可以亲手操作的例子，但如果你把它理解透彻了，你也就理解了一般的\trDerivedFunctor。$\mathrm{Tor}$ 将有助于证明关于\trFlatness{}的若干事实，我们将在 \S24.2 中讨论平坦性。$\mathrm{Tor}$ 是「torsion」（\trTorsion）的缩写（见\mbox{注记~24.2.1}）。

如果你从未见过这个概念，你或许只须记住它的性质即可。不过我反正会把一切都证明一遍——出奇地容易。

$\mathrm{Tor}$ 背后的想法如下。每当见到一个\trRightExact{}函子，我们总是希望它是一条\trLongExactSequence{}的末端。非正式地说，给定\trShortExactSequence
\[
\tag{23.1.0.1}\label{eq:tor-ses}
0 \longrightarrow N' \longrightarrow N \longrightarrow N'' \longrightarrow 0,
\]
我们希望 $M\otimes_A N' \to M\otimes_A N \to M\otimes_A N'' \to 0$ 能延拓成一条\trLongExactSequence
\begin{equation*}\tag{23.1.0.2}\label{eq:tor-les}
\begin{tikzcd}[row sep=2.9em, column sep=2.2em]
\cdots \ar[r] & \mathrm{Tor}^A_i(M,N') \ar[r] & \mathrm{Tor}^A_i(M,N) \ar[r] & \mathrm{Tor}^A_i(M,N'') \ar[r] & \cdots \\
 \ar[r] & \mathrm{Tor}^A_1(M,N') \ar[r] & \mathrm{Tor}^A_1(M,N) \ar[r] & \mathrm{Tor}^A_1(M,N'') & \\
 \ar[r] & M\otimes_A N' \ar[r] & M\otimes_A N \ar[r] & M\otimes_A N'' \ar[r] & 0.
\end{tikzcd}
\end{equation*}

更精确地说，我们所期望的是\emph{\trCovariantFunctor}{} $\mathrm{Tor}^A_i(M,\cdot)$：从 $A$-\trModule{}到 $A$-\trModule{}（其\trFunctoriality{}给出 \eqref{eq:tor-les} 中 $2/3$ 的态射），满足 $\mathrm{Tor}^A_0(M,N)\equiv M\otimes_A N$，并带有自然的\emph{\trConnectingHomomorphism}
\[
\delta\colon \mathrm{Tor}^A_{i+1}(M,N'') \longrightarrow \mathrm{Tor}^A_i(M,N')
\]
使得每个\trShortExactSequence{} \eqref{eq:tor-ses} 都给出\trLongExactSequence{} \eqref{eq:tor-les}。（「自然的」意思是：给定\trShortExactSequence{}之间的一个态射，你会写下的那个涉及 $\delta$-态射的自然方块必须交换。）

事实证明，把这件事做成并不太难，而这也将顺带为\trDerivedFunctor{}提供动机。我们现在来定义 $\mathrm{Tor}^A_i(M,N)$。

取 $N$ 的任意一个由自由模构成的分解 $\mathscr{R}$：
\[
\cdots \longrightarrow A^{\oplus n_2} \longrightarrow A^{\oplus n_1} \longrightarrow A^{\oplus n_0} \longrightarrow N \longrightarrow 0.
\]

更精确地说，从右往左构造这个分解。先选取 $N$ 作为 $A$-\trModule{}的一组\trGenerator，这便给出 $A^{\oplus n_0}\to N\to 0$；然后选取\trKernel{}的\trGenerator，依此类推。注意我们并不要求诸 $n_i$ 有限（尽管若 $N$ 是\trFinitelyGenerated{}的\trModule{}而 $A$ \trNoether，本来是可以取成有限的）。截断这个分解：把最后一项 $N$ 剥掉（把 $\to N\to 0$ 换成 $\to 0$）。然后与 $M$ 作张量（这并\emph{不}保持正合性）。注意 $M\otimes A^{\oplus n_1}=M^{\oplus n_1}$，因为与 $M$ 作张量和任意\trDirectSum{}交换（\mbox{练习~1.2.M}）。把 $\mathrm{Tor}^A_i(M,N)_{\mathscr{R}}$ 定义为这个\trComplexN{}在第 $i$ 级（$i\ge 0$）处的\trHomology。下标 $\mathscr{R}$ 提醒我们：这个构造依赖于所选取的分解，尽管我们很快就会看到，它其实与 $\mathscr{R}$ 无关。

我们来作几条快速的观察。
\begin{itemize}
\item $\mathrm{Tor}^A_0(M,N)_{\mathscr{R}}\cong M\otimes_A N$，\trCanonically。理由：由于作张量是\trRightExact{}的，而 $A^{\oplus n_1}\to A^{\oplus n_0}\to N\to 0$ 正合，故 $M^{\oplus n_1}\to M^{\oplus n_0}\to M\otimes_A N\to 0$ 正合，从而截断后的\trComplexN{} $M^{\oplus n_1}\to M^{\oplus n_0}\to 0$ 之\trHomology{}为 $M\otimes_A N$。
\item 若 $M\otimes\cdot$ 是\trExact{}的（即 $M$ \trFlat，\S1.5.13），则对一切 $i>0$ 有 $\mathrm{Tor}^A_i(M,N)_{\mathscr{R}}=0$。（这刻画了\trFlatness；见\mbox{练习~23.1.C}。）
\end{itemize}

现在给定两个\trModule{} $N$ 与 $N'$，以及它们各自的分解 $\mathscr{R}$ 与 $\mathscr{R}'$，我们可以把任意态射 $N\to N'$ 「提升」为这两个分解之间的一个态射：
\[
\begin{tikzcd}[row sep=2.9em, column sep=2em]
\cdots \ar[r] & A^{\oplus n_i} \ar[r] \ar[d] & \cdots \ar[r] & A^{\oplus n_1} \ar[r] \ar[d] & A^{\oplus n_0} \ar[r] \ar[d] & N \ar[r] \ar[d] & 0 \\
\cdots \ar[r] & A^{\oplus n'_i} \ar[r] & \cdots \ar[r] & A^{\oplus n'_1} \ar[r] & A^{\oplus n'_0} \ar[r] & N' \ar[r] & 0
\end{tikzcd}
\]
我们对 $i$ 归纳地做这件事。这里用到了 $A^{\oplus n_i}$ 的自由性：设 $a_1,\dots,a_{n_i}$ 是 $A^{\oplus n_i}$ 的\trGenerator，为把映射 $b\colon A^{\oplus n_i}\to A^{\oplus n'_{i-1}}$（它是微分 $A^{\oplus n_i}\to A^{\oplus n_{i-1}}$ 与先前构造好的 $A^{\oplus n_{i-1}}\to A^{\oplus n'_{i-1}}$ 的复合）提升为 $c\colon A^{\oplus n_i}\to A^{\oplus n'_i}$，我们随意地把诸 $b(a_j)$ 从 $A^{\oplus n'_{i-1}}$ 提升到 $A^{\oplus n'_i}$，并宣布这就是 $c(a_j)$。（给在意这类事情的人一个警示：我们在这里用到了选择公理。）

把这一组提升的选取记作 $\mathscr{R}\to\mathscr{R}'$。现在把两个\trComplexes{}都截断（去掉 $N\to N'$ 这一列），再与 $M$ 作张量。\trComplexes{}之间的映射诱导出\trHomology{}上的映射（\mbox{练习~1.5.D}），于是我们就描述了如下一些映射（先验地依赖于 $\mathscr{R}\to\mathscr{R}'$）：
\[
\mathrm{Tor}^A_i(M,N)_{\mathscr{R}} \longrightarrow \mathrm{Tor}^A_i(M,N')_{\mathscr{R}'}.
\]

回想（\S1.5.7）：若存在一列映射 $w\colon C_i\to C'_{i+1}$ 使得 $f-g=dw+wd$，则称\trComplexes{}的两个映射 $f\colon C_\bullet\to C'_\bullet$ 与 $g\colon C_\bullet\to C'_\bullet$ 是\trHomotopic。我们见过，两个同伦的映射在\trHomology{}上给出同一个映射（\mbox{练习~1.5.E}）。

\textbf{23.1.A.（关键练习）}\quad 证明：任意两个提升 $\mathscr{R}\to\mathscr{R}'$ 都是\trHomotopic。

现在我们把这些观察拢在一起。（请务必把它们彻底消化！）

（1）我们得到一个 $A$-\trModule{}的映射 $\mathrm{Tor}^A_i(M,N)_{\mathscr{R}}\to\mathrm{Tor}^A_i(M,N')_{\mathscr{R}'}$，它无关于提升 $\mathscr{R}\to\mathscr{R}'$ 的选取。

（2）从而，对 $A$-\trModule{} $N$ 的任意两个分解 $\mathscr{R}$ 与 $\mathscr{R}'$，我们得到一个\emph{\trCanonical}同构 $\mathrm{Tor}^A_i(M,N)_{\mathscr{R}}\overset{\sim}{\longrightarrow}\mathrm{Tor}^A_i(M,N)_{\mathscr{R}'}$。缘由如下。选取提升 $\mathscr{R}\to\mathscr{R}'$ 与 $\mathscr{R}'\to\mathscr{R}$。复合 $\mathscr{R}\to\mathscr{R}'\to\mathscr{R}$ 与恒等\trHomotopy（因为它是恒等映射 $N\to N$ 的一个提升）。于是，若 $f_{\mathscr{R}\to\mathscr{R}'}\colon\mathrm{Tor}^A_i(M,N)_{\mathscr{R}}\to\mathrm{Tor}^A_i(M,N)_{\mathscr{R}'}$ 是由 $\mathscr{R}\to\mathscr{R}'$ 诱导的映射，且类似地 $f_{\mathscr{R}'\to\mathscr{R}}$ 是由 $\mathscr{R}'\to\mathscr{R}$ 诱导的映射，则 $f_{\mathscr{R}'\to\mathscr{R}}\circ f_{\mathscr{R}\to\mathscr{R}'}$ 是恒等，同理 $f_{\mathscr{R}\to\mathscr{R}'}\circ f_{\mathscr{R}'\to\mathscr{R}}$ 亦然。

（3）从而 $\mathrm{Tor}^A_i(M,\cdot)$ 无关于分解的选取。它是一个\trCovariantFunctor{} $\mathit{Mod}_A\to\mathit{Mod}_A$。

\textbf{23.1.1.}\ \emph{\trhRemark。}\quad 注意：若 $N$ 是一个自由\trModule，则对一切 $M$ 与一切 $i>0$ 都有 $\mathrm{Tor}^A_i(M,N)=0$，因为 $N$ 具有平凡的分解 $0\to N\to N\to 0$（它是「自己的分解」）。

最后，我们得到\trLongExactSequence：

\textbf{23.1.2. \trhProposition。}\ ——\ \emph{对任意\trShortExactSequence{} \eqref{eq:tor-ses}，我们得到诸 $\mathrm{Tor}$ 的\trLongExactSequence{} \eqref{eq:tor-les}。}

\emph{证明。}\quad 给定一个\trShortExactSequence{} \eqref{eq:tor-ses}，选取 $N'$ 与 $N''$ 的分解。然后用它们按下述方式得到 $N$ 的一个分解。
\begin{equation*}\tag{23.1.2.1}\label{eq:tor-horseshoe}
\begin{tikzcd}[row sep=2.4em, column sep=2.2em]
 & 0 \ar[d] & 0 \ar[d] & 0 \ar[d] & \\
\cdots \ar[r] & A^{\oplus n'_i} \ar[r] \ar[d] & A^{\oplus n'_0} \ar[r] \ar[d] & N' \ar[r] \ar[d] & 0 \\
\cdots \ar[r] & A^{\oplus(n'_i+n''_i)} \ar[r] \ar[d] & A^{\oplus(n'_0+n''_0)} \ar[r] \ar[d] & N \ar[r] \ar[d] & 0 \\
\cdots \ar[r] & A^{\oplus n''_i} \ar[r] \ar[d] & A^{\oplus n''_0} \ar[r] \ar[d] & N'' \ar[r] \ar[d] & 0 \\
 & 0 & 0 & 0 &
\end{tikzcd}
\end{equation*}
映射 $A^{\oplus(n'_{i+1}+n''_{i+1})}\to A^{\oplus(n'_i+n''_i)}$ 是复合 $A^{\oplus n'_{i+1}}\to A^{\oplus n'_i}\hookrightarrow A^{\oplus(n'_i+n''_i)}$ 连同 $A^{\oplus n''_{i+1}}\to A^{\oplus n''_i}$ 到 $A^{\oplus(n'_i+n''_i)}$ 的一个提升一起给出的，这就保证中行是一个\trComplexN。

\textbf{23.1.B.（练习）}\quad 验证：可以选取从 $A^{\oplus n''_{i+1}}\to A^{\oplus n''_i}$ 到 $A^{\oplus(n'_i+n''_i)}$ 的这样一个提升。

于是，利用上同调的\trLongExactSequence（\mbox{定理~1.5.8}）以及顶行与底行正合这两个事实，\eqref{eq:tor-horseshoe} 的中行是\emph{正合}的（而不只是一个\trComplexN）。从而中行是一个分解，而 \eqref{eq:tor-horseshoe} 是分解的一个\trShortExactSequence。（这种构造有时被称作\emph{马蹄构造}（horseshoe construction），因为补全中行的样子，就像补全马蹄的中间部分。）注意到「$N$ 列」以外的各列都是「直和正合列」，而中行里的诸水平映射是「块上三角」的，这或许会有帮助。

然后截断（去掉最右一列 $0\to N'\to N\to N''\to 0$），与 $M$ 作张量（得到\trComplexes{}的一个\trShortExactSequence），再取上同调，便得到所要的\trLongExactSequence。\hfill$\square$

\textbf{23.1.C.（练习：$\mathrm{Tor}_1$ 平坦性判别法）}\quad 证明：以下条件关于 $A$-\trModule{} $M$ 等价。

（i）$M$ 是\trFlat{}的。

（ii）对一切 $i>0$ 与一切 $A$-\trModule{} $N$，$\mathrm{Tor}^A_i(M,N)=0$。

（iii）对一切 $A$-\trModule{} $N$，$\mathrm{Tor}^A_1(M,N)=0$。

\textbf{23.1.3.}\ \emph{\trhCaveat。}\quad 鉴于自由模显然是\trFlat{}的，你可能会以为\mbox{练习~23.1.C} 蕴含\mbox{注记~23.1.1}。如果我们知道 $\mathrm{Tor}^A_i(M,N)\cong\mathrm{Tor}^A_i(N,M)$，这的确会随之而来——当 $i=0$ 时这是清楚的（因为 $\otimes$ 是对称的），但当 $i>0$ 时，关于 $\mathrm{Tor}_i$ 的这件事我们要到\mbox{练习~23.3.A} 才会知道。

\textbf{23.1.D.（练习）}\quad 证明：上文构造的\trConnectingHomomorphism{} $\delta$ 无关于一切选择（诸分解等）。请以尽可能少的不快把这件事做出来。（可能的提示：给定用于构造 \eqref{eq:tor-horseshoe} 的两组选择，构造一个映射——一张三维的图——把 \eqref{eq:tor-horseshoe} 的一个版本连到另一个版本。）

\textbf{23.1.E.（不重要练习）}\quad 证明：$\mathrm{Tor}^A_i(M,\cdot)$ 是一个\trAdditiveFunctor（\mbox{定义~1.5.1}）。（我们后面不会用到它，你大可跳过。）

这样，我们就打好了 $\mathrm{Tor}$ 的地基。




\section{一般导出函子}

\textbf{23.2.1. \trProjectiveObj{}分解。}\quad 在上述 $\mathrm{Tor}$ 的构造中，我们关于自由模所用到的东西其实很少——事实上，我们只用到「自由模是\textbf{\trProjectiveObj}的」这一点，也就是说：它们是那样的\trModule{} $P$，对任何\trSurjectiveMap{} $M\twoheadrightarrow N$，都可以把任意态射 $P\to N$ 提升为 $P\to M$：
\begin{equation*}\tag{23.2.1.1}\label{eq:proj-lifting}
\begin{tikzcd}[row sep=2.4em, column sep=3.2em]
 & P \ar[dl, "\text{exists}"', dashed] \ar[dr] & \\
M \ar[rr, two heads] & & N
\end{tikzcd}
\end{equation*}
（正如 \S23.1 中所指出的，这用到了选择公理。）等价地，$\mathrm{Hom}(P,\cdot)$ 是一个\trExact{}函子（回想 $\mathrm{Hom}(Q,\cdot)$ 对任何 $Q$ 总是\trLeftExact{}的）。更一般地，同一个想法给出了「\textbf{任何\trAbelianCategory{}中的\trProjectiveObj{}对象}」的定义。于是，把我们的整个论证中的自由模统统换成\trProjectiveObj{}模，便可：(i) 取 $N$ 的任意一个\trProjectiveObj{}分解来计算 $\mathrm{Tor}^A_i(M,N)$；(ii) 对任何\trProjectiveObj{} $A$-\trModule{} $N$，有 $\mathrm{Tor}^A_i(M,N)=0$。

\textbf{23.2.A.（练习）}\quad 证明：一个对象 $P$ 是\trProjectiveObj{}的，当且仅当每个\trShortExactSequence{} $0\to Q\to R\to P\to 0$ 都分裂。从而证明：$A$-\trModule{} $P$ 是\trProjectiveObj{}的，当且仅当 $P$ 是某个自由\trModule{}的直和项。

\textbf{23.2.B.（练习）}\quad 证明：\trProjectiveObj{}模都是\trFlat{}的。提示：\mbox{练习~23.2.A}。若你想动用\mbox{练习~23.1.C}，请务必当心；见\mbox{警示~23.1.3}。（事实上，局部环上有限生成的\trProjectiveObj{}模甚至是自由的；见\mbox{注记~24.3.6}。）

\textbf{23.2.2. \trhDefinition{}：导出函子。}\quad 上述描述本是土法，却立刻可以剧烈地推广。我们所用到的无非是：$M\otimes_A\cdot$ 是一个\trRightExact{}函子，并且对任何 $A$-\trModule{} $N$，我们都能找到一个从\trProjectiveObj{}\trModule{}出发的\trSurjectiveMap{} $P\twoheadrightarrow N$。一般地，若 $F$ 是从 $A$-\trModule{}的范畴到任意\trAbelianCategory{}的\emph{任何}一个\trRightExact{}\trCovariantFunctor，这个构造都会定义出一列函子 $\mathrm{L}_iF$，使得 $\mathrm{L}_0F=F$，而诸 $\mathrm{L}_iF$ 给出一条\trLongExactSequence。我们还可以让它更一般。我们说一个\trAbelianCategory{}\textbf{有足够多\trProjectiveObj{}对象}（has enough projectives），如果对其中任何对象 $N$，都存在一个从\trProjectiveObj{}对象出发而满射到 $N$ 的态射。那么，若 $F$ 是从任何带足够多\trProjectiveObj{}对象的\trAbelianCategory{}到任何\trAbelianCategory{}的\trRightExact{}\trCovariantFunctor，我们就可以定义 $F$ 的\textbf{左导出函子}，记作 $\mathrm{L}_iF$（$i\ge 0$）。你应当重读 \S23.1，看到我们通篇只用到「我们有\trProjectiveObj{}分解（按 \eqref{eq:proj-lifting} 反复提升映射）」这一事实，以及「$F$ 把乘积送到乘积」（这是函子加性的推论，见\mbox{注记~1.5.2}），便可证明 $F$ 作用到 \eqref{eq:tor-horseshoe} 上保持各列的正合性。

\textbf{23.2.C.（练习）}\quad \textbf{内射对象}（injective object）这一概念，正是\trAbelianCategory{}中\trProjectiveObj{}对象概念的对偶。
（a）请把内射对象的定义精确地陈述出来。
（b）为以下各类函子定义导出函子：(i) 协变\trLeftExact{}函子（这些称作\textbf{右导出函子}）；(ii) 反变\trLeftExact{}函子（也称作\textbf{右导出函子}）；(iii) 反变\trRightExact{}函子（这些称作\textbf{左导出函子}），并把「该范畴有足够多内射对象或\trProjectiveObj{}对象」的必要假设显式写明。

\textbf{23.2.3.}\ \emph{\trhNotation。}\quad 回想\mbox{定义~23.2.2}：若 $F$ 是一个\trRightExact{}函子，则其（左）导出函子记作 $\mathrm{L}_iF$（$i\ge 0$，其中 $\mathrm{L}_0F=F$）。类似地，若 $F$ 是一个\trLeftExact{}函子，则其（右）导出函子记作 $\mathrm{R}^iF$。这个 $i$ 是上标，用以表明该\trLongExactSequence{}「随 $i$ 递升」。

\textbf{23.2.4. $\mathrm{Ext}$ 函子。}\quad

\textbf{23.2.D.（容易练习（及定义）：$A$-\trModule{}的 $\mathrm{Ext}$ 函子，第一版）}\quad $\mathrm{Hom}(\cdot,N)$ 是 $\mathit{Mod}_A$ 上一个反变\trLeftExact{}函子，而 $\mathit{Mod}_A$ 有足够多\trProjectiveObj{}对象。定义 $\mathrm{Ext}^i_A(M,N)$ 为 $\mathrm{Hom}(\cdot,N)$ 的第 $i$ 个右导出函子作用在 $M$ 上。请陈述诸 $\mathrm{Ext}$ 模的相应\trLongExactSequence。

\textbf{23.2.E.（容易练习（及定义）：$A$-\trModule{}的 $\mathrm{Ext}$ 函子，第二版）}\quad 范畴 $\mathit{Mod}_A$ 有足够多内射对象（见 \S23.2.5）。由于 $\mathrm{Hom}(M,\cdot)$ 是 $\mathit{Mod}_A$ 上一个协变\trLeftExact{}函子，定义 $\mathrm{Ext}^i_A(M,N)$ 为 $\mathrm{Hom}(M,\cdot)$ 的右导出函子作用在 $N$ 上。请陈述诸 $\mathrm{Ext}$ 模的相应\trLongExactSequence。

看起来前两道练习出了问题：我们把 $\mathrm{Ext}^i(M,N)$ 定义了两遍，还得到了两条不同的\trLongExactSequence！所幸，$\mathrm{Ext}$ 的这两个定义是一致的（见\mbox{练习~23.3.B}），而对 $\mathrm{Ext}$ 来说，两条长正合列总比一条好。

\textbf{23.2.F.（容易练习）}\quad 用\mbox{练习~23.2.D} 中 $\mathrm{Ext}$ 的定义证明：若 $A$ 是\trNoether 环，而 $M$ 与 $N$ 是\trFinitelyGenerated{}的 $A$-\trModule，则 $\mathrm{Ext}^i_A(M,N)$ 是一个\trFinitelyGenerated{}的 $A$-\trModule。

（$\mathrm{Ext}$-函子（就\trSheaf{}而言）将在\trSerreDuality{}的证明中扮演关键角色；见 \S29.2。）

\textbf{23.2.5.}\ \hard\ \textbf{$A$-\trModule{}的范畴有足够多内射对象。}\quad 我们将需要「$\mathit{Mod}_A$ 有足够多内射对象」这一事实，但其证明细节日后不会再露面，所以你尽管跳过这段讨论。

\textbf{23.2.G.（练习）}\quad 设 $Q$ 是一个 $A$-\trModule，它使得对每个\trIdeal{} $I\subset A$，每个\trHomomorphism{} $I\to Q$ 都能延拓为 $A\to Q$。证明 $Q$ 是一个内射 $A$-\trModule。提示：设 $N\subset M$ 是 $A$-\trModule{}的一个包含，而我们给定 $\beta\colon N\to Q$。我们希望证明 $\beta$ 能延拓为 $M\to Q$。用\trZornsLemma 证明：在那些满足 $N\subset N'\subset M$ 且 $\beta$ 能延拓为 $N'\to Q$ 的 $A$-\trModule{} $N'$ 之中，存在一个极大的。若这个 $N'$ 不是 $M$，请给出 $\beta$ 到 $N'+Am$ 上的一个延拓，其中 $m\in M\setminus N'$，从而导出矛盾。

\textbf{23.2.H.（容易练习（用到\trZornsLemma））}\quad 证明：一个 $\mathbb{Z}$-\trModule（即\trAbelianGroups）$Q$ 是内射的，当且仅当它是\trDivisibleGroup（即：对每个 $q\in Q$ 与每个 $n\in\mathbb{Z}^{\neq 0}$，存在 $q'\in Q$ 使得 $nq'=q$）。从而证明：内射 $\mathbb{Z}$-\trModule{}的任何\trQuotient{}也是内射的。

\textbf{23.2.I.（练习）}\quad 证明：$\mathbb{Z}$-\trModule{}的范畴 $\mathit{Mod}_{\mathbb{Z}}=\mathit{Ab}$ 有足够多内射对象。（提示：若 $M$ 是一个 $\mathbb{Z}$-\trModule，把它写成自由 $\mathbb{Z}$-\trModule{} $F$ 对某个 $K$ 取商的结果。证明 $M$ 含于可除群 $(F\otimes_{\mathbb{Z}}\mathbb{Q})/K$ 之中。）

\textbf{23.2.J.（练习）}\quad 设 $Q$ 是一个内射 $\mathbb{Z}$-\trModule，而 $A$ 是一个\trRing。证明 $\mathrm{Hom}_{\mathbb{Z}}(A,Q)$ 是一个内射 $A$-\trModule。提示：先描述 $\mathrm{Hom}_{\mathbb{Z}}(A,Q)$ 上的 $A$-\trModule{}结构。你只会用到「$\mathbb{Z}$ 是一个环」以及「$A$ 是该环上的一个\trAlgebra」这两件事。

\textbf{23.2.K.（练习）}\quad 证明 $\mathit{Mod}_A$ 有足够多内射对象。提示：设 $M$ 是一个 $A$-\trModule。由\mbox{练习~23.2.I}，我们可以找到一个\emph{$\mathbb{Z}$-模}的包含 $M\hookrightarrow Q$，其中 $Q$ 是一个内射 \emph{$\mathbb{Z}$-模}。请描述 $A$-\trModule{}包含的一个序列
\[
M \hookrightarrow \mathrm{Hom}_{\mathbb{Z}}(A,M) \hookrightarrow \mathrm{Hom}_{\mathbb{Z}}(A,Q).
\]
（$\mathrm{Hom}_{\mathbb{Z}}(A,M)$ 上的 $A$-\trModule{}结构经由左变元 $A$ 上的 $A$-作用给出，而非经由右变元 $M$ 上的 $A$-作用给出。）右端那一项由前一道的\mbox{练习~23.2.J} 是内射的。

\textbf{23.2.6.}\ \hard\ \textbf{泛\trDeltaFunctor。}\quad （这段讨论在头一遍阅读时最好跳过；你应当直接前往 \S23.3。在 \S29.2.2 之前我们不会用到 $\delta$-函子这个概念。）现在我们来描述\trDerivedFunctor{}的一个更一般的变体，在第 29 章讨论\trSerreDuality{}时你可能就会用到它。泛\trDeltaFunctor{}这个概念的优点在于：即便在 $\mathscr{A}$ 没有足够多内射对象的情形，我们也能应用它。

把\trDerivedFunctor{}的关键性质抽象出来，我们定义如下资料：一个（上同调性的）\trDeltaFunctor，它从一个\trAbelianCategory{} $\mathscr{A}$ 出发，落到另一个\trAbelianCategory{} $\mathscr{B}$。一个\textbf{\trDeltaFunctor}是一族加性函子 $\mathrm{T}^i\colon\mathscr{A}\to\mathscr{B}$（其中 $i<0$ 时取 $\mathrm{T}^i$ 为 $0$），连同对每个\trShortExactSequence
\[
\tag{23.2.6.1}\label{eq:delta-ses}
0 \longrightarrow A' \longrightarrow A \longrightarrow A'' \longrightarrow 0
\]
所给出的态射 $\delta^i\colon\mathrm{T}^i(A'')\to\mathrm{T}^{i+1}(A')$，满足两条性质：

（i）（\emph{短正合列给出长正合列}）对每个\trShortExactSequence{} \eqref{eq:delta-ses}，序列
\[
\begin{tikzcd}[column sep=1.7em]
\cdots \ar[r] & \mathrm{T}^{i-1}(A'') \ar[r, "\delta^{i-1}"] & \mathrm{T}^i(A') \ar[r] & \mathrm{T}^i(A) \ar[r] & \mathrm{T}^i(A'') \ar[r, "\delta^i"] & \mathrm{T}^{i+1}(A') \ar[r] & \cdots
\end{tikzcd}
\]
（其中未标出的映射来自诸 $\mathrm{T}^i$ 的协变性）是正合的。特别地，$\mathrm{T}^0$ 是\trLeftExact{}的。

（ii）（(i) 的函子性）对于 $\mathscr{A}$ 中\trShortExactSequence{}之间的每个态射
\[
\begin{tikzcd}[row sep=2.2em, column sep=2.2em]
0 \ar[r] & \mathfrak{a}' \ar[r] \ar[d] & \mathfrak{a} \ar[r] \ar[d] & \mathfrak{a}'' \ar[r] \ar[d] & 0 \\
0 \ar[r] & A' \ar[r] & A \ar[r] & A'' \ar[r] & 0
\end{tikzcd}
\]
（其中诸方块交换），诸 $\delta^i$ 给出交换图
\[
\begin{tikzcd}[row sep=2.6em, column sep=3.4em]
\mathrm{T}^i(\mathfrak{a}'') \ar[r, "\delta^i"] \ar[d] & \mathrm{T}^{i+1}(\mathfrak{a}') \ar[d] \\
\mathrm{T}^i(A'') \ar[r, "\delta^i"] & \mathrm{T}^{i+1}(A')
\end{tikzcd}
\]
（其中诸竖直箭头来自 $\mathrm{T}^i$ 与 $\mathrm{T}^{i+1}$ 的协变性）。

导出函子上同调显然是一个\trDeltaFunctor{}的例子；拟紧分离概形上\trSheaf{}的\trCechCohomology{}是另一个。（你愿意的话，可以把这两句话说精确。）

\textbf{23.2.7.}\ \emph{\trhRemark／\trhDefinition。}\quad 更精确地说，我们应当把上者称作\textbf{协变\trDeltaFunctor}。类似地还有\textbf{反变\trDeltaFunctor}的概念，其定义你当能自行猜出。

\textbf{23.2.L.（练习）}\quad 想清楚 $\mathscr{A}\to\mathscr{B}$ 的\textbf{$\delta$-函子态射}（morphism of $\delta$-functors）的恰当定义。（随后应当清楚：从 $\mathscr{A}$ 到 $\mathscr{B}$ 的诸 $\delta$-函子构成一个范畴。）

\textbf{23.2.8.}\ \emph{\trhDefinition。}\quad 一个（上同调性的）\trDeltaFunctor{} $(\mathrm{T}^i,\delta^i)$ 是\textbf{泛}的，如果对任何其他\trDeltaFunctor{} $(\mathrm{T}'^i,\delta'^i)$，以及函子的任何\trNaturalTransformation{} $\alpha\colon\mathrm{T}^0\to\mathrm{T}'^0$，都存在唯一的\trDeltaFunctor{}态射 $(\mathrm{T}^i,\delta^i)\to(\mathrm{T}'^i,\delta'^i)$ 延拓 $\alpha$。由泛性质的废话（以及\mbox{练习~23.2.L}），给定任何协变\trLeftExact{}函子 $F\colon\mathscr{A}\to\mathscr{B}$，至多存在一个延拓 $F$ 的泛\trDeltaFunctor{} $(\mathrm{T}^i,\delta^i)$（即带有自然同构 $\mathrm{T}^0\overset{\sim}{\longleftarrow}F$ 的那一个）。关于泛\trDeltaFunctor{}的关键事实是下述结果。

\textbf{23.2.9. \trhTheorem。}\ ——\ \emph{设 $(\mathrm{T}^i,\delta^i)$ 是从 $\mathscr{A}$ 到 $\mathscr{B}$ 的一个协变\trDeltaFunctor，且对一切 $A\in\mathscr{A}$，都存在一个单态射 $A\to J$ 使得对一切 $i>0$ 有 $\mathrm{T}^i(J)=0$。则 $(\mathrm{T}^i,\delta^i)$ 是泛的。}

\textbf{23.2.M.}\ \extrahard\ \textbf{（练习）}\quad 证明\mbox{定理~23.2.9}。部分提示：受下方\mbox{推论~23.2.11} 的启发，沿循我们对\trDerivedFunctor{}的讨论。更好的提示（因为这道练习很难）：沿循 [Weib, Exercise 2.4.5] 中的提示，或沿循 [Lan, Ch.~XX, Thm.~7.1] 的证明。

\textbf{23.2.10.}\ \emph{\trhRemark。}\quad 一个加性函子 $F\colon\mathscr{A}\to\mathscr{B}$ 称作\textbf{\trEffaceable{}的}（effaceable），如果对每个 $A\in\mathscr{A}$，都存在一个单态射 $A\to J$ 使得 $\mathrm{F}(J)=0$。\mbox{定理~23.2.9} 的假设可以削弱为只要求每个 $\mathrm{T}^i$（$i>0$）都\trEffaceable{}的，而改证这较弱的一版，你也定会欣然从命。（确实，[Weib, Exercise 2.4.5]、[Lan, Ch.~XX, Thm.~7.1] 以及原始文献 [Gr1, II.2.2.1] 处理的就是这种情形。）我们陈述 \mbox{定理~23.2.9} 时取其简形式，因为我们只会用到这个版本。

\textbf{23.2.11. \trhCorollary。}\ ——\ \emph{若 $\mathscr{A}$ 有足够多内射对象，而 $F$ 是一个协变\trLeftExact{}函子 $\mathscr{A}\to\mathscr{B}$，则诸 $\mathrm{R}^iF$（连同伴随它们的诸 $\delta^i$）构成一个泛\trDeltaFunctor。}

\emph{证明。}\quad $\mathscr{A}$ 的每个元素都容纳入一个到内射对象的单态射；这不过是「有足够多内射对象」的定义（\mbox{练习~23.2.C}）。内射对象 $I$ 的高阶导出函子总是 $0$：只须取 $I$ 「自己的」内射分解来计算其高阶导出函子。\hfill$\square$



\section{导出函子与谱序列}

不少有用的事实都可以用\trSpectralSequence{}轻松证明。做完这些练习，你就会在对类似情形动用谱序列论证时不再有任何恐惧，因为你将意识到：它们全都是同一回事。

在你阅读本节之前，你应当先读一下讲\trSpectralSequence{}的 \S1.6。

\textbf{23.3.1. $\mathrm{Tor}$ 的对称性。}\quad

\textbf{23.3.A.（练习：$\mathrm{Tor}$ 的对称性）}\quad 证明：存在一个同构 $\mathrm{Tor}^A_1(M,N)\overset{\sim}{\longrightarrow}\mathrm{Tor}^A_1(N,M)$。（提示：取 $M$ 的一个自由分解与 $N$ 的一个自由分解。把二者的「乘积」以某种方式做成一个\trDoubleComplex。把这个明显的\trSpectralSequence{}的两个朝向都用上，看看你会得到什么。）

顺带说一句相关的事：

\textbf{23.3.B.（练习）}\quad 证明：\mbox{练习~23.2.D} 与\mbox{练习~23.2.E} 中给出的 $\mathrm{Ext}^1(M,N)$ 的两个定义是一致的。

\textbf{23.3.2. 导出函子可以用零调分解来计算。}\quad 设 $F\colon\mathscr{A}\to\mathscr{B}$ 是\trAbelianCategory{}之间的一个\trRightExact{}加性函子，并且 $\mathscr{A}$ 有足够多\trProjectiveObj{}对象。我们称 $A\in\mathscr{A}$ 是 \textbf{$F$-\trAcyclic}的（或当 $F$ 在上下文中清楚时简称\textbf{\trAcyclic}的），如果对 $i>0$ 有 $\mathrm{L}_iF\,A=0$。在\mbox{练习~23.3.D} 中我们将看到，$\mathscr{A}$ 的对象 $B$ 的导出函子 $\mathrm{L}_iF\,B$ 可以用「零调分解」来计算。我们先用一个有用的构造搭好舞台。

\textbf{23.3.3.}\ \emph{为一个正合序列构造「\trProjectiveObj{}分解」。}\quad 设 $\cdots\to E_2\to E_1\to E_0\to 0$ 是某个带足够多\trProjectiveObj{}对象的\trAbelianCategory{}中的一个正合序列。我们来解释如何归纳地造出一个由\trProjectiveObj{}对象构成的\trDoubleComplex
\[
\begin{tikzcd}[row sep=2.2em, column sep=1.9em]
 & \vdots \ar[d] & \vdots \ar[d] & \vdots \ar[d] & \\
\cdots \ar[r] & P_{2,1} \ar[r] \ar[d] & P_{1,1} \ar[r] \ar[d] & P_{0,1} \ar[r] \ar[d] & 0 \\
\cdots \ar[r] & P_{2,0} \ar[r] \ar[d] & P_{1,0} \ar[r] \ar[d] & P_{0,0} \ar[r] \ar[d] & 0 \\
\cdots \ar[r] & E_2 \ar[r] \ar[d] & E_1 \ar[r] \ar[d] & E_0 \ar[r] \ar[d] & 0 \\
 & 0 & 0 & 0 &
\end{tikzcd}
\]
使得各行与各列都正合。假设你已经造好了这个\trComplexN{}的一部分，正打算造出 $P_{m,n}$ 这一项：
\begin{equation*}\tag{23.3.3.1}\label{eq:pp-lifting}
\begin{tikzcd}[row sep=2.2em, column sep=2.4em]
? \ar[r] & P_{m-1,n} \ar[r] \ar[d] & P_{m-2,n} \ar[d] \\
P_{m,n-1} \ar[r] \ar[d] & P_{m-1,n-1} \ar[r] \ar[d] & P_{m-2,n-1} \\
P_{m,n-2} \ar[r] & P_{m-1,n-2} &
\end{tikzcd}
\end{equation*}
出于一些很快就会明朗的理由，我们假设：对任意 $b\in P_{m-1,n-1}$，只要它在 $P_{m-1,n-2}$ 与 $P_{m-2,n-1}$ 中的像都为零，就存在 $c\in P_{m,n-1}$，它在 $P_{m-1,n-1}$ 中的像为 $b$，而在 $P_{m,n-2}$ 中的像为零；并且对称地，存在 $c'\in P_{m-1,n}$，它在 $P_{m-1,n-1}$ 中的像为 $b$，而在 $P_{m-2,n}$ 中的像为零。

考虑
\[
K:=\ker\bigl(P_{m,n-1}\oplus P_{m-1,n}\to P_{m,n-2}\oplus P_{m-1,n-1}\oplus P_{m-2,n}\bigr).
\]
若 $K$ 满射到 $\ker(P_{m-1,n}\to P_{m-2,n})$ 与 $\ker(P_{m,n-1}\to P_{m,n-2})$ 二者之上，那么我们就可以取从\trProjectiveObj{}对象 $P\twoheadrightarrow K$ 出发的任意满态射，然后取 $P_{m,n}$ 为 $P$（配备其到 $P_{m,n-1}$ 的映射，以及其到 $P_{m-1,n}$ 的映射之负）。这样一来，我们就同时保证了 $P_{m-1,n}$ 处的「水平正合性」、$P_{m,n-1}$ 处的「竖直正合性」，以及 (23.3.3.1) 中方块的交换性。

我们现在用「追图」（这是合法的，见 \S1.5.5）来验证这两个想要的满态射。设 $a\in\ker(P_{m-1,n}\to P_{m-2,n})$，并设 $b$ 是它在 $P_{m-1,n-1}$ 中的像。我们希望找到 $c\in\ker(P_{m,n-1}\to P_{m,n-2})$，使得 $c$ 在 $P_{m-1,n-1}$ 中映到 $b$，而在 $P_{m,n-2}$ 中映到 $0$。（那么 $(-c,a)\in P_{m,n-1}\oplus P_{m-1,n}$ 就会是我们想要的、在 $P_{m-1,n}$ 中映到 $a$ 的那个 $K$ 的元素。）但这样的 $c$ 依我们的假设是存在的\footnote{原书此处作「exits」，当为「exists」之误，径改。——译者注}！对称地，我们便得到满态射 $K\twoheadrightarrow\ker(P_{m,n-1}\to P_{m,n-2})$。

最后一项任务，是保证这些假设在我们随后建造\trDoubleComplex{}的各阶段依然成立。我们需要保证：对任意 $b'\in P_{m-1,n}$，只要它映到 $(0,0)\in P_{m-1,n-1}\oplus P_{m-2,n}$，就有 $P_{m,n}$ 的一个元素映到 $(0,b')\in P_{m,n-1}\oplus P_{m-1,n}$。并且我们必须保证把 $m$ 与 $n$ 的角色互换后的类似陈述。我们取 $Q'=\ker(P_{m-1,n}\to P_{m-1,n-1}\oplus P_{m-2,n})$，取 $Q''=\ker(P_{m,n-1}\to P_{m,n-2}\oplus P_{m-1,n-1})$，并取（从\trProjectiveObj{}对象出发的）满态射 $P'\twoheadrightarrow Q'$ 与 $P''\twoheadrightarrow Q''$。然后取 $P_{m,n}$ 为 $P'\oplus P''$，其中从 $P'$ 出发的映射如下：$P'$ 到 $P_{m-1,n}$ 的映射是上面描述的那个映射（经由 $Q'$）；$P'$ 到 $P_{m,n-1}$ 的映射为零；而对 $P''$ 则相反（到 $P_{m,n-1}$ 的映射为零，而到 $P_{m-1,n}$ 的映射如上经由 $Q''$ 所蕴含）。和项 $P'$ 保证我们的第一条想要成立的假设，和项 $P''$ 保证第二条。

\textbf{23.3.C.（练习）}\quad 验证：上述构造确实给出一个正合序列的\trProjectiveObj{}分解。你是在哪里用到了「序列 $E_\bullet$ 正合」这一点的？

现在，让我们来应用它。

\textbf{23.3.D.（练习）}\quad 证明：你可以用\textbf{零调分解}（acyclic resolutions，而不仅仅是\trProjectiveObj{}分解）来计算 $\mathscr{A}$ 的对象 $B$ 的导出函子，也就是说，取一个由\trAcyclic{}对象构成的分解
\[
\tag{23.3.3.2}\label{eq:acyclic-resolution}
\cdots \longrightarrow A_2 \longrightarrow A_1 \longrightarrow A_0 \longrightarrow B \longrightarrow 0
\]
截断，作用 $F$，然后取\trHomology。从而，$\mathrm{Tor}_i(M,N)$ 可以用 $M$ 或 $N$ 的一个平坦分解来计算。提示：如上所述，在正合序列 \eqref{eq:acyclic-resolution} 「上方」造一个由\trProjectiveObj{}对象构成的\trDoubleComplex。去掉最底下一行与最右边非零的一列，然后把 $F$ 作用上去，得到一个新的\trDoubleComplex。用\trSpectralSequence{}论证证明：(i) 该\trDoubleComplex{}的\trHomology{}等于 $\mathrm{L}_iF(B)$；(ii) 该\trDoubleComplex{}的\trHomology{}与练习题面所给的构造相一致。如果这让你晕头转向，请阅读下方关于\trCartan--\trEilenberg{}分解的更多内容。

\textbf{23.3.4. \trGrothendieck 函子复合谱序列。}\quad 设 $\mathscr{A}$、$\mathscr{B}$ 与 $\mathscr{C}$ 是\trAbelianCategory，$F\colon\mathscr{A}\to\mathscr{B}$ 与 $G\colon\mathscr{B}\to\mathscr{C}$ 是\trLeftExact{}加性协变函子，并且 $\mathscr{A}$ 与 $\mathscr{B}$ 有足够多内射对象。于是，$F$、$G$ 与 $G\circ F$ 的右导出函子都存在。一个合理的问题是：它们之间有何关系？

\textbf{23.3.5. \trhTheorem（\trGrothendieck 函子复合谱序列）。}\ ——\ \emph{设 $F\colon\mathscr{A}\to\mathscr{B}$ 与 $G\colon\mathscr{B}\to\mathscr{C}$ 是\trLeftExact{}加性协变函子，且 $\mathscr{A}$ 与 $\mathscr{B}$ 有足够多内射对象。再设 $F$ 把 $\mathscr{A}$ 的内射对象送到 $\mathscr{B}$ 的 $G$-\trAcyclic{}对象。则对每个 $X\in\mathscr{A}$，都存在一个\trSpectralSequence，满足 $\Eto{p,q}{2}=\mathrm{R}^q\mathrm{G}(\mathrm{R}^p\mathrm{F}(X))$，抵达 $\mathrm{R}^\bullet(\mathrm{G}\circ\mathrm{F})(X)$。}

我们很快就会看到勒雷\trSpectralSequence{}作为一个应用（\mbox{定理~23.4.5}）。

从\mbox{定理~23.3.5} 的证明中还有一些东西值得榨取。例如，尽管谱序列的 $E_0$ 页将依赖于若干（内射分解的）选择，$E_2$ 页却将与选择无关。对我们的应用而言，我们并不需要这种加细。

建立\mbox{定理~23.3.5} 是要下一番功夫的，因此其证明在头一遍阅读时或许最好跳过。

\textbf{23.3.6.}\ \hard\ \textbf{证明\mbox{定理~23.3.5}。}\quad 在给出证明（\S23.3.8）之前，我们先做些准备工作来为之提供动机。为了讨论作用在 $X$ 上的导出函子，我们选取 $X$ 的一个内射分解：
\[
0 \longrightarrow X \longrightarrow I^0 \longrightarrow I^1 \longrightarrow \cdots.
\]
为了计算导出函子 $\mathrm{R}^p\mathrm{F}(X)$，我们把 $F$ 作用到内射分解 $I^\bullet$ 上：
\[
0 \longrightarrow F(I^0) \longrightarrow F(I^1) \longrightarrow F(I^2) \longrightarrow \cdots.
\]
注意，依\mbox{定理~23.3.5} 的假设，$F(I^\bullet)$ 是 $G$-\trAcyclic{}的。如果我们顺着自己的鼻子走，或许会对上面这个\trComplexN{} $F(I^\bullet)$ 的诸项取「同时」内射分解 $I^{\bullet,\bullet}$（\S23.3.3 的「对偶」——注意只有各列被要求正合），然后作用 $G$，并审视所得到的\trDoubleComplex：
\begin{equation*}\tag{23.3.6.1}\label{eq:ce-double}
\begin{tikzcd}[row sep=2.2em, column sep=2em]
 & \vdots & \vdots & \vdots & \\
0 \ar[r] & G(I^{0,2}) \ar[r] \ar[u] & G(I^{1,2}) \ar[r] \ar[u] & G(I^{2,2}) \ar[r] \ar[u] & \cdots \\
0 \ar[r] & G(I^{0,1}) \ar[r] \ar[u] & G(I^{1,1}) \ar[r] \ar[u] & G(I^{2,1}) \ar[r] \ar[u] & \cdots \\
0 \ar[r] & G(I^{0,0}) \ar[r] \ar[u] & G(I^{1,0}) \ar[r] \ar[u] & G(I^{2,0}) \ar[r] \ar[u] & \cdots \\
 & 0 \ar[u] & 0 \ar[u] & 0. \ar[u] &
\end{tikzcd}
\end{equation*}

\textbf{23.3.E.（练习）}\quad 考虑带上行\trOrientation{}、以 \eqref{eq:ce-double} 为 $E_0$ \trPage{}的\trSpectralSequence。证明：$E_2^{p,q}$ 在 $q=0$ 时为 $\mathrm{R}^p(\mathrm{G}\circ\mathrm{F})(X)$，否则为 $0$。

我们现在看到了\mbox{定理~23.3.5} 结论中一半的项；我们已经走到了半途。为了完成证明，我们还想考虑另一个带右行\trOrientation{}的谱序列，但我们需要对 \eqref{eq:ce-double} 知道得更多；我们将更细心地把它造出来。

\textbf{23.3.7.}\ \emph{\trCartan--\trEilenberg{}分解。}\quad 设 $\cdots\to C^{p-1}\to C^p\to C^{p+1}\to\cdots$ 是\trAbelianCategory{} $\mathscr{B}$ 中的一个\trComplexN。我们将为 $C^\bullet$ 造一个内射分解
\begin{equation*}\tag{23.3.7.1}\label{eq:ce-resolution}
\begin{tikzcd}[row sep=2em, column sep=1.8em]
 & \vdots & \vdots & \vdots & \\
0 \ar[r] & I^{0,2} \ar[r] \ar[u] & I^{1,2} \ar[r] \ar[u] & I^{2,2} \ar[r] \ar[u] & \cdots \\
0 \ar[r] & I^{0,1} \ar[r] \ar[u] & I^{1,1} \ar[r] \ar[u] & I^{2,1} \ar[r] \ar[u] & \cdots \\
0 \ar[r] & I^{0,0} \ar[r] \ar[u] & I^{1,0} \ar[r] \ar[u] & I^{2,0} \ar[r] \ar[u] & \cdots \\
0 \ar[r] & C^0 \ar[r] \ar[u] & C^1 \ar[r] \ar[u] & C^2 \ar[r] \ar[u] & \cdots \\
 & 0 \ar[u] & 0 \ar[u] & 0 \ar[u] &
\end{tikzcd}
\end{equation*}
并满足一些进一步的性质。

我们先为\trComplexN{}上的若干函数定义一些记号。
\begin{itemize}
\item 置 $\mathrm{Z}^p(K^\bullet)$ 为\trComplexN{} $K^\bullet$ 的第 $p$ 个微分之\trKernel。
\item 置 $\mathrm{B}^{p+1}(K^\bullet)$ 为\trComplexN{} $K^\bullet$ 的第 $p$ 个微分之像。（上标如此选取，使得 $\mathrm{B}^{p+1}(K^\bullet)\subset\mathrm{Z}^{p+1}$\footnote{原书此处作「$\mathrm{B}^{p+1}(K^\bullet)\subset\mathrm{K}^{p+1}$」，当为「$\mathrm{Z}^{p+1}$」之误（$\mathrm{B}^{p+1}$ 含于 $\mathrm{K}^{p+1}$ 系平凡事实，此括注意在说明上标的选取使诸边缘落入诸循环），径改。——译者注}。）
\item 照例，置 $\mathrm{H}^p(K^\bullet)$ 为\trComplexN{} $K^\bullet$ 在第 $p$ 步处的\trCohomology。
\end{itemize}
对每个 $p$，我们都有\trComplexes
\begin{equation*}\tag{23.3.7.2}\label{eq:ce-three}
\begin{tikzcd}[row sep=1.9em, column sep=1.6em]
 & \vdots & \vdots & \vdots & \\
 & \mathrm{Z}^p(I^{\bullet,1}) \ar[u] & \mathrm{B}^p(I^{\bullet,1}) \ar[u] & \mathrm{H}^p(I^{\bullet,1}) \ar[u] & \\
 & \mathrm{Z}^p(I^{\bullet,0}) \ar[u] & \mathrm{B}^p(I^{\bullet,0}) \ar[u] & \mathrm{H}^p(I^{\bullet,0}) \ar[u] & \\
 & \mathrm{Z}^p(C^\bullet) \ar[u] & \mathrm{B}^p(C^\bullet) \ar[u] & \mathrm{H}^p(C^\bullet) \ar[u] & \\
 & 0 \ar[u] & 0 \ar[u] & 0. \ar[u] &
\end{tikzcd}
\end{equation*}
我们将这样构造 \eqref{eq:ce-resolution}，使得 (23.3.7.2) 的三个\trComplexes{}都是内射分解（就其首个非零项而言）。我们先把 $\mathrm{B}^p(C^\bullet)$ 的内射分解 $\mathrm{B}^{p,\ast}$ 与 $\mathrm{H}^p(C^\bullet)$ 的内射分解 $\mathrm{H}^{p,\ast}$ 选取好；它们最终将成为 (23.3.7.2) 的最后两行。

\textbf{23.3.F.（练习）}\quad 请描述 $\mathrm{Z}^p(C^\bullet)$ 的一个内射分解 $\mathrm{Z}^{p,\ast}$（即 (23.3.7.2) 的第一行），使得下图成为\trComplexes{}的一个\trShortExactSequence。
\begin{equation*}\tag{23.3.7.3}\label{eq:ce-bzh-diagram}
\begin{tikzcd}[row sep=2em, column sep=1.8em]
 & \vdots & \vdots & \vdots & \\
0 \ar[r] & \mathrm{B}^{p,1} \ar[r] \ar[u] & \mathrm{Z}^{p,1} \ar[r] \ar[u] & \mathrm{H}^{p,1} \ar[r] \ar[u] & 0 \\
0 \ar[r] & \mathrm{B}^{p,0} \ar[r] \ar[u] & \mathrm{Z}^{p,0} \ar[r] \ar[u] & \mathrm{H}^{p,0} \ar[r] \ar[u] & 0 \\
0 \ar[r] & \mathrm{B}^p(C^\bullet) \ar[r] \ar[u] & \mathrm{Z}^p(C^\bullet) \ar[r] \ar[u] & \mathrm{H}^p(C^\bullet) \ar[r] \ar[u] & 0 \\
 & 0 \ar[u] & 0 \ar[u] & 0 \ar[u] &
\end{tikzcd}
\end{equation*}
提示：「对偶」的问题已在 \eqref{eq:tor-horseshoe} 中以「马蹄构造」解决了。

\textbf{23.3.G.（练习）}\quad 请描述 $C^p$ 的一个内射分解 $\mathrm{I}^{p,\ast}$，使得下图成为\trComplexes{}的一个\trShortExactSequence。
\begin{equation*}\tag{23.3.7.4}\label{eq:ce-zib-diagram}
\begin{tikzcd}[row sep=2em, column sep=1.8em]
 & \vdots & \vdots & \vdots & \\
0 \ar[r] & \mathrm{Z}^{p,1} \ar[r] \ar[u] & \mathrm{I}^{p,1} \ar[r] \ar[u] & \mathrm{B}^{p+1,1} \ar[r] \ar[u] & 0 \\
0 \ar[r] & \mathrm{Z}^{p,0} \ar[r] \ar[u] & \mathrm{I}^{p,0} \ar[r] \ar[u] & \mathrm{B}^{p+1,0} \ar[r] \ar[u] & 0 \\
0 \ar[r] & \mathrm{Z}^p(C^\bullet) \ar[r] \ar[u] & C^p \ar[r] \ar[u] & \mathrm{B}^{p+1}(C^\bullet) \ar[r] \ar[u] & 0 \\
 & 0 \ar[u] & 0 \ar[u] & 0 \ar[u] &
\end{tikzcd}
\end{equation*}
（前一题的提示在这里再度适用。我们指出，(23.3.7.3) 与 (23.3.7.4) 最底下的非零行曾出现于 \mbox{(1.5.6.3)} 中。）

\textbf{23.3.H.（练习／定义）}\quad 构造 $C^\bullet$ 的一个内射分解 \eqref{eq:ce-resolution}，使得 $\mathrm{Z}^{p,\ast}=\mathrm{Z}^p(\mathrm{I}^{\bullet,\ast})$、$\mathrm{B}^{p,\ast}=\mathrm{B}^p(\mathrm{I}^{\bullet,\ast})$、$\mathrm{H}^{p,\ast}=\mathrm{H}^p(\mathrm{I}^{\bullet,\ast})$，从而 (23.3.7.2) 的三个\trComplexes{}都是内射分解。这称作 $C^\bullet$ 的一个\textbf{\trCartan--\trEilenberg{}分解}。

\textbf{23.3.8.}\ \emph{\trGrothendieck{}谱序列的证明，\mbox{定理~23.3.5}。}\quad 我们从插叙\trCartan--\trEilenberg{}分解之前停下的地方继续。选取 $X$ 的一个内射分解 $I^\bullet$。把函子 $F$ 作用上去，然后取 $\mathrm{F}I^\bullet$ 的一个\trCartan--\trEilenberg{}分解 $\mathrm{I}^{\bullet,\bullet}$，再作用 $G$，便得到 \eqref{eq:ce-double}。

\mbox{练习~23.3.E} 描述了把 \eqref{eq:ce-double} 作为 $E_0$ \trPage{}、取上行\trOrientation{}的谱序列会发生什么。于是我们现在考虑右行\trOrientation。

从我们对\trCartan--\trEilenberg{}分解的构造，我们得到了诸内射分解 \eqref{eq:ce-three}，以及\trShortExactSequence
\[
\tag{23.3.8.1}\label{eq:ce-bzh}
0 \longrightarrow \mathrm{B}^p(\mathrm{I}^{\bullet,q}) \longrightarrow \mathrm{Z}^p(\mathrm{I}^{\bullet,q}) \longrightarrow \mathrm{H}^p(\mathrm{I}^{\bullet,q}) \longrightarrow 0,
\]
\[
\tag{23.3.8.2}\label{eq:ce-zib}
0 \longrightarrow \mathrm{Z}^p(\mathrm{I}^{\bullet,q}) \longrightarrow \mathrm{I}^{p,q} \longrightarrow \mathrm{B}^{p+1}(\mathrm{I}^{\bullet,q}) \longrightarrow 0
\]
——它们由\emph{内射}对象构成（来自 (23.3.7.3) 与 (23.3.7.4) 的诸列）。这意味着二者都是\emph{分裂}正合列（中间一项可以被表达为外侧两项的直和），于是作用 $G$ 之后，两条正合列依然正合。

把\trLeftExact{}函子 $G$ 作用到
\[
0 \longrightarrow \mathrm{Z}^p(\mathrm{I}^{\bullet,q}) \longrightarrow \mathrm{I}^{p,q} \longrightarrow \mathrm{B}^{p+1}(\mathrm{I}^{\bullet,q}),
\]
我们发现 $\mathrm{GZ}^p(\mathrm{I}^{\bullet,q})=\ker(\mathrm{GI}^{p,q}\to\mathrm{GI}^{p+1,q})$。但这个\trKernel{}正是 $\mathrm{Z}^p(\mathrm{GI}^{\bullet,q})$ 的\emph{定义}，于是我们得到一个诱导同构 $\mathrm{GZ}^p(\mathrm{I}^{\bullet,q})=\mathrm{Z}^p(\mathrm{GI}^{\bullet,q})$（「$G$ 与 $\mathrm{Z}^p$ 交换」）。由 (23.3.8.2) 的正合性，在作用 $G$ 后我们看到 $\mathrm{GB}^{p+1}(\mathrm{I}^{\bullet,q})=\mathrm{B}^{p+1}(\mathrm{GI}^{\bullet,q})$（二者都是 $\coker(\mathrm{GZ}^p(\mathrm{I}^{\bullet,q})\to\mathrm{GI}^{p,q})$）。由 (23.3.8.1) 的正合性，在作用 $G$ 后我们看到 $\mathrm{GH}^p(\mathrm{I}^{\bullet,q})=\mathrm{H}^p(\mathrm{GI}^{\bullet,q})$（二者都是 $\coker(\mathrm{GB}^p(\mathrm{I}^{\bullet,q})\to\mathrm{GZ}^p(\mathrm{I}^{\bullet,q}))$——于是「$G$ 与 $\mathrm{H}^p$ 交换」）。

我们回到带右行\trOrientation{}、以 \eqref{eq:ce-double} 为 $E_0$ \trPage{}的谱序列。在右行方向上取\trCohomology，我们发现 $E_1^{p,q}=\mathrm{H}^p(\mathrm{GI}^{\bullet,q})=\mathrm{GH}^p(\mathrm{I}^{\bullet,q})$（因为 $G$ 与 $\mathrm{H}^p$ 交换）。现在 $\mathrm{H}^p(\mathrm{I}^{\bullet,q})$ 是 $(\mathrm{R}^p\mathrm{F})(X)$ 的一个内射分解（\eqref{eq:ce-three} 中的最后一个分解）。于是，当我们用竖直箭头计算 $E_2$ 时，便得到 $E_2^{p,q}=\mathrm{R}^q\mathrm{G}(\mathrm{R}^p\mathrm{F}(X))$。

你现在应当自行验证：这一点（与\mbox{练习~23.3.E} 相结合）便结束了\mbox{定理~23.3.5} 的证明。\hfill$\square$




\section{$\OO_X$ 模的导出函子上同调}

我们想把导出函子的这套机器用起来，为概形 $X$ 上的\trQCSheaf{}定义上同调。与其在范畴 $\mathit{QCoh}_X$ 中工作，出于若干理由，在更大的范畴 $\mathit{Mod}_{\OO_X}$ 中工作会更简单（但参看\mbox{不重要注记~23.5.7}）。

\textbf{23.4.1. \trhTheorem。}\ ——\ \emph{设 $(X,\OO_X)$ 是一个\trRingedSpace。则 $\OO_X$-\trModule{}的范畴 $\mathit{Mod}_{\OO_X}$ 有足够多内射对象。}

作为一项副产品（对别人比对我们更有用），取 $\OO_X=\mathbb{Z}$，我们看到：固定拓扑空间上阿贝尔群层的范畴有足够多内射对象。

\textbf{23.4.2.}\ \emph{证明。}\quad 我们沿循 Godement 的做法 [GrS, pp.~27--28]，用一系列练习来证明\mbox{定理~23.4.1}。设 $\mathscr{F}$ 是一个 $\OO_X$-\trModule。我们将展示一个单射 $\mathscr{F}\hookrightarrow\mathscr{Q}'$，其中 $\mathscr{Q}'$ 是某个内射 $\OO_X$-\trModule。对每个 $\mathfrak{p}\in X$，选取一个包含 $\mathscr{F}_\mathfrak{p}\hookrightarrow Q_\mathfrak{p}$，其中 $Q_\mathfrak{p}$ 是一个内射 $\OO_{X,\mathfrak{p}}$-\trModule（这是可能的，因为 $\OO_{X,\mathfrak{p}}$-\trModule{}的范畴有足够多内射对象，\mbox{练习~23.2.K}）。

\textbf{23.4.A.（练习）}\quad 证明：\trSkyscraperSheaf{} $\mathscr{Q}_\mathfrak{p}:=i_{\mathfrak{p},\ast}Q_\mathfrak{p}$（带有位于点 $\mathfrak{p}\in X$ 处的\trModule{} $Q_\mathfrak{p}$）是一个内射 $\OO_X$-\trModule。（你可以用\trAbstractNonsense{}漂亮地做到这一点——动用\mbox{练习~23.5.B}——但动手去做也同样快。）

\textbf{23.4.B.（容易练习）}\quad 证明：\trAbelianCategory{}中内射对象的（可能是无限的）乘积仍是内射的。

由前两道练习，$\mathscr{Q}':=\prod_{\mathfrak{p}\in X}\mathscr{Q}_\mathfrak{p}$ 是一个内射 $\OO_X$-\trModule。

\textbf{23.4.C.（容易练习）}\quad 通过审视茎，证明自然映射 $\mathscr{F}\to\mathscr{Q}'$ 是一个单射。

这样，\mbox{定理~23.4.1} 的证明便告完成。\hfill$\square$

我们现在可以给出好几个定义了。

\textbf{23.4.3.}\ \emph{\trhDefinition。}\quad 若 $(X,\OO_X)$ 是一个\trRingedSpace，而 $\mathscr{F}$ 是一个 $\OO_X$-\trModule，则定义 $\mathrm{H}^i(X,\mathscr{F})$ 为 $\mathrm{R}^i\Gamma(X,\mathscr{F})$。此外，若 $\pi\colon(X,\OO_X)\to(Y,\OO_Y)$ 是\trRingedSpace{}之间的一个映射，则我们有\textbf{\trHigherPushforward}（或\textbf{导出前推}）$\mathrm{R}^i\pi_\ast\colon\mathit{Mod}_{\OO_X}\to\mathit{Mod}_{\OO_Y}$。

这些概念我们早先在特殊情形中定义过：第 18 章中，对拟紧分离概形上的\trQCSheaf{}（就 $\mathrm{H}^i$ 而言），或对概形的拟紧分离态射（就 $\mathrm{R}^i\pi_\ast$ 而言）。我们很快将在 \S23.5 中证明，这些较早的定义与\mbox{定义~23.4.3} 相一致。因此，导出函子的定义比我们的切赫定义适用面广泛得多。但值得指出的是，我们将要用到的几乎所有东西都可以从切赫定义中得出。一个值得注意的例外是\trLeray{}\trSpectralSequence，我们现在就来讨论它。

\textbf{23.4.4. \trLeray{}谱序列。}\quad

\textbf{23.4.5. \trhTheorem（\trLeray{}谱序列）。}\ ——\ \emph{设 $\pi\colon(X,\OO_X)\to(Y,\OO_Y)$ 是\trRingedSpace{}之间的一个\trMorphism。对任何 $\OO_X$-\trModule{} $\mathscr{F}$，存在一个谱序列，其 $E_2$ 项由 $\mathrm{H}^q(Y,\mathrm{R}^p\pi_\ast\mathscr{F})$ 给出，抵达 $\mathrm{H}^{p+q}(X,\mathscr{F})$。}

这是\trGrothendieck 函子复合谱序列（\mbox{定理~23.3.5}）的一个直接推论——只须证明内射 $\OO$-\trModule{}的前推是\trAcyclic{}的 $\OO$-\trModule。我们现在就来证明这件事。

\textbf{23.4.6.}\ \emph{\trhDefinition。}\quad 我们给出一个本身也很重要的中间定义。拓扑空间上的层 $\mathscr{F}$ 称作\textbf{\trFlasque}的（有时也称作 flabby），如果一切限制映射都是满射，也就是说，如果 $\mathrm{res}_{\mathfrak{U}\subset\mathfrak{V}}\colon\mathscr{F}(\mathfrak{V})\to\mathscr{F}(\mathfrak{U})$ 对一切 $\mathfrak{U}\subset\mathfrak{V}$ 都是满射。

\textbf{23.4.D.（练习）}\quad 设 $(X,\OO_X)$ 是一个\trRingedSpace。
（a）证明：若
\[
\tag{23.4.6.1}\label{eq:flasque-ses}
0 \longrightarrow \mathscr{F}' \longrightarrow \mathscr{F} \longrightarrow \mathscr{F}'' \longrightarrow 0
\]
是 $\OO_X$-\trModule{}的一个正合序列，而 $\mathscr{F}'$ 是\trFlasque{}的，则 (23.4.6.1) 在任意开集 $\mathfrak{U}$ 的截面上正合，即 $0\to\mathscr{F}'(\mathfrak{U})\to\mathscr{F}(\mathfrak{U})\to\mathscr{F}''(\mathfrak{U})\to 0$
是正合的。
（b）给定正合序列 \eqref{eq:flasque-ses}，若 $\mathscr{F}'$ 是\trFlasque{}的，证明：$\mathscr{F}$ 是\trFlasque{}的当且仅当 $\mathscr{F}''$ 是\trFlasque{}的。

\textbf{23.4.E.（容易练习：松软层的前推仍是松软层）}\quad
（a）设 $\pi\colon X\to Y$ 是拓扑空间之间的一个连续映射，而 $\mathscr{F}$ 是 $X$ 上的一个由集合构成的\trFlasque{}层。证明 $\pi_\ast\mathscr{F}$ 是 $Y$ 上的一个\trFlasque{}层。
（b）设 $\pi\colon(X,\OO_X)\to(Y,\OO_Y)$ 是\trRingedSpace{}之间的一个\trMorphism，而 $\mathscr{F}$ 是一个\trFlasque{}的 $\OO_X$-\trModule。证明 $\pi_\ast\mathscr{F}$ 是一个\trFlasque{}的 $\OO_Y$-\trModule。

\textbf{23.4.7.}\ \emph{零扩张：逆像函子的一个偶尔扮演的左伴随。}\quad 为了证明内射层是\trFlasque{}的（\mbox{练习~23.4.H}），我们引入「\textbf{\trExtensionByZero}函子」。除了总是一个左伴随之外，当 $\pi$ 是开子集的包含时，$\pi^{-1}$ 有时还可以是一个右伴随。设 $i\colon\mathfrak{U}\hookrightarrow Y$ 是开子集到 $Y$ 中的一个包含。如下定义\textbf{\trExtensionByZero}函子 $i_!\colon\mathit{Mod}_{\OO_\mathfrak{U}}\to\mathit{Mod}_{\OO_Y}$。设 $\mathscr{F}$ 是一个 $\OO_\mathfrak{U}$-\trModule。对 $W\subset Y$，定义 $(i^{\mathrm{pre}}\mathscr{F})(W)=\mathscr{F}(W)$（若 $W\subset\mathfrak{U}$），否则为 $0$（限制映射取显然的那些）。这显然是一个预层 $\OO_Y$-\trModule。定义 $i_!$ 为 $(i^{\mathrm{pre}})^{\mathrm{sh}}$。注意 $i_!\mathscr{F}$ 是一个 $\OO_Y$-\trModule，并且这定义了一个函子。（符号「$!$」读作「shriek」（尖叫），大概是因为人们看到它就会尖叫。于是 $i_!$ 读作「i-lower-shriek」。）

\textbf{23.4.F.（练习）}\quad 证明：$i^{\mathrm{pre}}\mathscr{F}$ 满足恒等公理，但未必是一个层。（我们后面不会用到这一点，但它或许能让你看出这为什么叫「零扩张」。可能的例子来源：$\mathbb{R}$ 上的连续函数。）

\textbf{23.4.G.（练习）}\quad
（a）对 $\mathfrak{q}\in Y$，证明 $(i_!\mathscr{F})_\mathfrak{q}=\mathscr{F}_\mathfrak{q}$（若 $\mathfrak{q}\in\mathfrak{U}$），否则为 $0$。
（b）证明 $i_!$ 是一个\trExact{}函子。
（c）若 $\mathscr{G}$ 是一个 $\OO_Y$-\trModule，请描述一个包含 $i_!i^{-1}\mathscr{G}\hookrightarrow\mathscr{G}$。（一个有趣的注记，我们用不到：设 $Z$ 为 $\mathfrak{U}$ 的补集，并设 $j\colon Z\to Y$ 为自然包含。那么就存在一个\trShortExactSequence
\[
0 \longrightarrow i_!i^{-1}\mathscr{G} \longrightarrow \mathscr{G} \longrightarrow j_\ast j^{-1}\mathscr{G} \longrightarrow 0.
\]
最好是通过描述这些映射、然后在茎上检验正合性来验证这一点。）
（d）证明 $(i_!,i^{-1})$ 是一个伴随对，从而存在自然双射
\[
\mathrm{Hom}_{\OO_Y}(i_!\mathscr{F},\mathscr{G})\overset{\sim}{\longrightarrow}\mathrm{Hom}_{\OO_\mathfrak{U}}(\mathscr{F},\mathscr{G}|_\mathfrak{U})
\]
对任何 $\OO_\mathfrak{U}$-\trModule{} $\mathscr{F}$ 与 $\OO_Y$-\trModule{} $\mathscr{G}$。（特别地，$\mathscr{G}$ 在 $\mathfrak{U}$ 上的截面可以与 $\mathrm{Hom}_{\OO_Y}(i_!\OO_\mathfrak{U},\mathscr{G})$ 等同起来。）

我们准备好使用\trExtensionByZero{}函子这件工具了。

\textbf{23.4.H.（练习：内射层是松软层）}\quad 设 $(X,\OO_X)$ 是一个\trRingedSpace，而 $\mathscr{Q}$ 是一个内射 $\OO_X$-\trModule。证明 $\mathscr{Q}$ 是\trFlasque{}的。提示：若 $\mathfrak{U}\subset\mathfrak{V}\subset X$，请描述 $\OO_X$-\trModule{}的一个单射 $0\to(i_\mathfrak{U})_!\OO_\mathfrak{U}\to(i_\mathfrak{V})_!\OO_\mathfrak{V}$，其中 $i_\mathfrak{U}\colon\mathfrak{U}\hookrightarrow X$ 与 $i_\mathfrak{V}\colon\mathfrak{V}\hookrightarrow X$ 是显然的开嵌入。对它作用正合的反变函子 $\mathrm{Hom}(\cdot,\mathscr{Q})$。

\textbf{23.4.I.（练习：松软层蕴含 $\Gamma$-\trAcyclic）}\quad 设 $\mathscr{F}$ 是一个\trFlasque{}的 $\OO_X$-\trModule。证明 $\mathscr{F}$ 是 $\Gamma$-\trAcyclic{}的（即对 $i>0$ 有 $\mathrm{H}^i(X,\mathscr{F})=0$，\S23.3.2），做法如下。由于 $\mathit{Mod}_{\OO_X}$ 有足够多内射对象，选取 $\mathscr{F}$ 进入某个内射 $\mathscr{I}$ 的一个包含，并称其余核为 $\mathscr{G}$：
\[
0 \longrightarrow \mathscr{F} \longrightarrow \mathscr{I} \longrightarrow \mathscr{G} \longrightarrow 0.
\]
那么由\mbox{练习~23.4.H}，$\mathscr{I}$ 是\trFlasque{}的；于是由\mbox{练习~23.4.D}(b)，$\mathscr{G}$ 是\trFlasque{}的。取（导出函子）上同调的\trLongExactSequence，证明 $\mathrm{H}^1(X,\mathscr{F})=0$。你的论证对\emph{任何}\trFlasque{}层 $\mathscr{F}$ 都有效，故 $\mathrm{H}^1(X,\mathscr{G})=0$ 也成立。证明 $\mathrm{H}^2(X,\mathscr{F})=0$。把这一点变成一次归纳。

于是，若 $\pi\colon X\to Y$ 是\trRingedSpace{}之间的一个\trMorphism，而 $\mathscr{I}$ 是一个内射 $\OO_X$-\trModule，则 $\mathscr{I}$ 是\trFlasque{}的（\mbox{练习~23.4.H}），于是 $\pi_\ast\mathscr{I}$ 是\trFlasque{}的（\mbox{练习~23.4.E}(b)），于是 $\pi_\ast\mathscr{I}$ 对函子 $\Gamma$ 而言是\trAcyclic{}的（\mbox{练习~23.4.I}），这就完成了\trLeray{}谱序列（\mbox{定理~23.4.5}）的证明。\hfill$\square$

\textbf{23.4.J.（练习）}\quad 把\trLeray{}谱序列（\mbox{定理~23.4.5}）推广到处理
\[
(X,\OO_X) \longrightarrow^{\pi} (Y,\OO_Y) \longrightarrow^{\rho} (Z,\OO_Z)
\]
的高阶前推的复合。

\textbf{23.4.8.}\ \extrahard\ \textbf{\trRingedSpace{}上的 $\OO_X$-\trModule{}范畴未必有足够多射影对象。}\quad 与\mbox{定理~23.4.1} 相对照，\trRingedSpace{}上的 $\OO_X$-\trModule{}的范畴未必有足够多射影对象。例如，设 $X$ 为带扎里斯基拓扑的 $\mathbb{P}^1_k$（事实上，我们对 $X$ 需要知道的非常少——只需要知道它不是一个 \trAlexandrov{}空间），但取 $\OO_X$ 为\trConstantSheaf{} $\mathbb{Z}$。我们将会看到 $\mathit{Mod}_{\OO_X}$——也就是说，$X$ 上阿贝尔群层的范畴——没有足够多射影对象。若 $\mathit{Mod}_{\OO_X}$ 有足够多射影对象，那么就会存在一个从射影层出发的满态射 $\psi\colon P\to\mathbb{Z}$。取定一个闭点 $\mathfrak{q}\in X$。我们将证明茎上的映射 $\psi_\mathfrak{q}\colon P_\mathfrak{q}\to\mathbb{Z}_\mathfrak{q}$ 是零映射，从而与 $\psi$ 的满性矛盾。对 $X$ 的每个开子集 $\mathfrak{U}$，把 $\mathbb{Z}_\mathfrak{U}$ 记作 $\mathfrak{U}$ 到 $X$ 的\trExtensionByZero{}（取值于 $\mathbb{Z}$ 的\trConstantSheaf）（\S23.4.7——$\mathbb{Z}_\mathfrak{U}(\mathfrak{V})=\mathbb{Z}$（若 $\mathfrak{V}\subset\mathfrak{U}$），否则 $\mathbb{Z}_\mathfrak{U}(\mathfrak{V})=0$）。对 $\mathfrak{q}$ 的每个开邻域 $\mathfrak{V}$，设 $W$ 是一个严格更小的开邻域。考虑满态射 $\mathbb{Z}_{X-\mathfrak{q}}\oplus\mathbb{Z}_W\to\mathbb{Z}$。由 $P$ 的射影性，满态射 $\psi$ 提升为 $P\to\mathbb{Z}_{X-\mathfrak{q}}\oplus\mathbb{Z}_W$。映射 $P(\mathfrak{V})\to\mathbb{Z}(\mathfrak{V})$ 穿过 $\mathbb{Z}_{X-\mathfrak{q}}(\mathfrak{V})\oplus\mathbb{Z}_W(\mathfrak{V})=0$ 分解，从而必是零映射。于是茎上的映射 $\psi_\mathfrak{q}\colon P_\mathfrak{q}\to\mathbb{Z}_\mathfrak{q}$ 也是零映射（你看出为什么了吗？），正如所愿。



\section{\trCechCohomology{}与导出函子上同调一致}

接下来我们要证明：在\trCechCohomology{}有定义之处，它与导出函子上同调一致。

\textbf{23.5.1. \trhTheorem。}\ ——\ \emph{设 $X$ 是一个拟紧分离概形，而 $\mathscr{F}$ 是一个\trQCSheaf。则 $\mathscr{F}$ 的\trCechCohomology{}与 $\mathscr{F}$ 的导出函子上同调一致。}

这句话还不如它本该有的那样精确。我们会想知道：这个同构对 $\mathscr{F}$ 是\trFunctoriality{}的，并且它尊重\trLongExactSequence（从而为\trCechCohomology{}定义的\trConnectingHomomorphism{}与为导出函子上同调定义者一致）。对\trHigherPushforward{}还有一个重要的推广。我们把这些问题留给本节末尾的 \S23.5.5。

出于好奇不妨一说：只要定义得当，那么在概形理论之外的一大圈情形中（若底拓扑空间仿紧），切赫上同调都与导出函子上同调一致——但并非永远如此（反例见 [Gr1, \S3.8]）。\mbox{注记~23.5.6} 也与此相关。

证明的核心想法（尽管带一个转折）是一种谱序列论证，风格与 \S23.3 中那些相同，并且对谱序列的每个朝向各用一条「上同调消失」式原料。

\textbf{(A)}~若 $(X,\OO_X)$ 是一个\trRingedSpace，$\mathscr{Q}$ 是一个内射 $\OO_X$-\trModule，而 $X=\cup_i\mathfrak{U}_i$ 是一个有限开覆盖，则对此覆盖而言，$\mathscr{Q}$ 没有第 $i$ 切赫上同调（$i>0$）。

\textbf{(B)}~若 $X$ 是一个\trAffineScheme，而 $\mathscr{F}$ 是 $X$ 上的一个\trQCSheaf，则 $(\mathrm{R}^i\Gamma)(\mathscr{F})=0$ 对 $i>0$。

翻译一下：\textbf{(A)} 说的是，导出函子上同调的积木块没有切赫上同调；而 \textbf{(B)} 说的是，切赫上同调的积木块没有导出函子上同调。

\textbf{23.5.A.（预备练习）}\quad 设 $(X,\OO_X)$ 是一个\trRingedSpace，$\mathscr{Q}$ 是一个内射 $\OO$-\trModule，而 $i\colon\mathfrak{U}\hookrightarrow X$ 是一个开子集。证明 $\mathscr{Q}|_\mathfrak{U}$ 在 $\mathfrak{U}$ 上是内射的。提示：动用「$i^{-1}$ 有一个\trExact{}的左伴随 $i_!$（\trExtensionByZero）」这一事实；见 \S23.4.7，以及下面两幅图。
\[
\begin{tikzcd}[row sep=2.6em, column sep=3em]
0 \ar[r] & \mathscr{A} \ar[r] \ar[d] & \mathscr{B} \ar[dl, dashed, "?"] & & 0 \ar[r] & i_!\mathscr{A} \ar[r] \ar[d] & i_!\mathscr{B} \ar[dl, dashed, "?"] \\
 & \mathscr{Q}|_\mathfrak{U} & & & & \mathscr{Q} &
\end{tikzcd}
\]

在做\mbox{练习~23.5.A} 的过程中，你将证明下面这条事实；它在\mbox{练习~29.4.A} 中还会再度用到。（你也可以用它来解\mbox{练习~23.4.A}。）

\textbf{23.5.B.（练习：正合函子的右伴随保持内射对象）}\quad 证明：若 $(F,G)$ 是\trAbelianCategory{}之间加性函子的一个\trAdjointPair，且 $F$ 是\trExact{}的，则 $G$ 把内射对象送到内射对象。

\textbf{23.5.2.}\ \emph{\mbox{定理~23.5.1} 的证明（假定 (A) 与 (B)）。}\quad 与 \S23.3 中所证的那些事实一样，我们只采用那条唯一合理的路线：选取 $\mathscr{F}$ 的一个内射分解 $0\to\mathscr{F}\to\mathscr{Q}_\bullet$ 与 $X$ 的一个切赫覆盖，把这两类信息在一个\trDoubleComplex{}中混起来，然后把它扔进我们的谱序列机器（\S1.6）。更精确地说：选取一个有限\trAffine{}开覆盖 $X=\cup_i\mathfrak{U}_i$ 与一个内射分解
\[
0 \longrightarrow \mathscr{F} \longrightarrow \mathscr{Q}_0 \longrightarrow \mathscr{Q}_1 \longrightarrow \cdots.
\]

考虑这个\trDoubleComplex
\begin{equation*}\tag{23.5.2.1}\label{eq:agree-double}
\begin{tikzcd}[row sep=2em, column sep=1.7em]
\vdots & \vdots & \vdots & \vdots & \\
0 \ar[r] \ar[u] & \oplus_i\mathscr{Q}_2(\mathfrak{U}_i) \ar[r] \ar[u] & \oplus_{i,j}\mathscr{Q}_2(\mathfrak{U}_{ij}) \ar[r] \ar[u] & \oplus_{i,j,k}\mathscr{Q}_2(\mathfrak{U}_{ijk}) \ar[r] \ar[u] & \cdots \\
0 \ar[r] \ar[u] & \oplus_i\mathscr{Q}_1(\mathfrak{U}_i) \ar[r] \ar[u] & \oplus_{i,j}\mathscr{Q}_1(\mathfrak{U}_{ij}) \ar[r] \ar[u] & \oplus_{i,j,k}\mathscr{Q}_1(\mathfrak{U}_{ijk}) \ar[r] \ar[u] & \cdots \\
0 \ar[r] \ar[u] & \oplus_i\mathscr{Q}_0(\mathfrak{U}_i) \ar[r] \ar[u] & \oplus_{i,j}\mathscr{Q}_0(\mathfrak{U}_{ij}) \ar[r] \ar[u] & \oplus_{i,j,k}\mathscr{Q}_0(\mathfrak{U}_{ijk}) \ar[r] \ar[u] & \cdots \\
0 \ar[r] \ar[u] & 0 \ar[r] \ar[u] & 0 \ar[r] \ar[u] & 0 \ar[r] \ar[u] & \cdots.
\end{tikzcd}
\end{equation*}

我们把它取作某个谱序列的 $E_0$ 项。先用右行滤过。由于内射 $\OO$-\trModule{}的高阶切赫上同调为 $0$（假设 \textbf{(A)}），除「第 $0$ 列」之外我们处处得到 $0$，而在第 $i$ 行的第 $0$ 列处我们得到 $\mathscr{Q}_i(X)$：
\[
\begin{tikzcd}[row sep=1.9em, column sep=1.4em]
 & \vdots & \vdots & \vdots & \vdots & \\
 & 0 \ar[u] & \mathscr{Q}_2(X) \ar[u] & 0 \ar[u] & 0 \ar[u] & \cdots \\
 & 0 \ar[u] & \mathscr{Q}_1(X) \ar[u] & 0 \ar[u] & 0 \ar[u] & \cdots \\
 & 0 \ar[u] & \mathscr{Q}_0(X) \ar[u] & 0 \ar[u] & 0 \ar[u] & \cdots \\
 & 0 \ar[u] & 0 \ar[u] & 0 \ar[u] & 0. \ar[u] &
\end{tikzcd}
\]

然后在竖直方向上取上同调，便在 $E_2$ 页上得到 $\mathscr{F}$ 在 $X$ 上的导出函子上同调：
\[
\begin{tikzcd}[row sep=2.2em, column sep=1.7em]
 & \vdots & \vdots & \vdots & \vdots & \\
 & 0 & \mathrm{R}^2\Gamma(X,\mathscr{F}) & 0 & 0 & \cdots \\
 & 0 & \mathrm{R}^1\Gamma(X,\mathscr{F}) & 0 & 0 & \cdots \\
 & 0 & \Gamma(X,\mathscr{F}) \ar[uul] & 0 \ar[uul] & 0 \ar[uul] & \cdots \\
 & & 0 \ar[uul] & 0 \ar[uul] & 0. \ar[uul] &
\end{tikzcd}
\]

然后我们回到 $E_0$ 页从头再来，这一次取上行定向（相应于先取向上的箭头）。由\mbox{命题~23.5.A}，$\mathscr{Q}_i|_{\mathfrak{U}_j}$ 在 $\mathfrak{U}_j$ 上是内射的，于是我们计算的是 $\mathscr{F}$ 在 $\mathfrak{U}_j$ 上的导出函子上同调。而 $\mathscr{F}$ 在 $\mathfrak{U}_j$ 上的高阶导出函子上同调为 $0$（假设 \textbf{(B)}），故除第 $0$ 行之外一切项都是 $0$。于是 $E_1$ 项为：
\begin{equation*}\tag{23.5.2.2}\label{eq:agree-e1}
\begin{tikzcd}[row sep=2.4em, column sep=2.4em]
0 \ar[r] & 0 \ar[r] & 0 \ar[r] & 0 \ar[r] & \cdots \\
0 \ar[r] & 0 \ar[r] & 0 \ar[r] & 0 \ar[r] & \cdots \\
0 \ar[r] & \oplus_i\Gamma(\mathfrak{U}_i,\mathscr{F}) \ar[r] & \oplus_{i,j}\Gamma(\mathfrak{U}_{ij},\mathscr{F}) \ar[r] & \oplus_{i,j,k}\Gamma(\mathfrak{U}_{ijk},\mathscr{F}) \ar[r] & \cdots \\
0 \ar[r] & 0 \ar[r] & 0 \ar[r] & 0. \ar[r] & \cdots.
\end{tikzcd}
\end{equation*}

第 $0$ 行恰好是 $\mathscr{F}$ 的切赫复形，于是该谱序列在 $E_2$ 项处收敛，给出切赫上同调。既然一个朝向给出导出函子上同调、另一个朝向给出切赫上同调，证明便告完成。\hfill$\square$

于是只剩下证明 \textbf{(A)} 与 \textbf{(B)}。

\textbf{23.5.3. 原料 (A)：内射对象没有切赫上同调。}\quad

\textbf{23.5.C.（练习）}\quad 设 $X=\cup_j\mathfrak{U}_j$ 是 $X$ 由开集构成的一个有限覆盖，而 $\mathscr{F}$ 是 $X$ 上的一个\trFlasque{}层（\mbox{定义~23.4.6}）。证明：$\mathscr{F}$ 相应于 $\cup_{j=1}^n\mathfrak{U}_j$ 的切赫复形在正次数处没有\trCohomology，也就是说，它在次数 $0$ 之外正合（在次数 $0$ 处它有\trCohomology{} $\mathscr{F}(X)$，依据的是层公理）。提示：对 $n$ 作归纳。考虑\trComplexes{}的\trShortExactSequence{} \mbox{(18.2.4.2)}（另见 \mbox{(18.2.3.1)}）。相应的\trLongExactSequence{}将立刻给出 $i>1$ 时所要的结果，而 $i=1$ 时则要用到松软性。

于是，\trFlasque{}层没有切赫上同调，因而内射 $\OO$-\trModule{}（\mbox{练习~23.4.H}）尤其也没有。对我们的代数几何应用而言，有这些就够了；但为了让你看看这套机器有多通用，我们给出一个有趣的应用。

\textbf{23.5.D.（不重要练习：借层上同调给出容斥原理的一条怪证）}\quad 容斥原理等价于下述命题：设 $X$ 是一个有限集，而 $\mathfrak{U}_i$（$1\le i\le n$）是覆盖 $X$ 的有限集。照例，对 $I\subset\{1,\dots,n\}$ 定义 $\mathfrak{U}_I=\cap_{i\in I}\mathfrak{U}_i$。那么
\[
|X|=\sum_i|\mathfrak{U}_i|-\sum_{|I|=2}|\mathfrak{U}_{I}|+\sum_{|I|=3}|\mathfrak{U}_{I}|-\sum_{|I|=4}|\mathfrak{U}_{I}|+\cdots.
\]
请给 $X$ 赋予\trDiscreteTopology，证明\trConstantSheaf{} $\mathbb{Q}$ 是\trFlasque{}的，考虑用覆盖 $\mathfrak{U}_i$ 计算 $\mathrm{H}^1(X,\mathbb{Q})$ 的切赫复形，并动用\mbox{练习~1.5.B}，由此证明上式。

\textbf{23.5.4.}\ \hard\ \textbf{原料 (B)：仿射概形上的拟凝聚层没有导出函子上同调。}\quad 我们对 $k$ 作归纳来证明下述命题。设 $X$ 是一个\trAffineScheme，而 $\mathscr{F}$ 是 $X$ 上的一个\trQCSheaf。则 $\mathrm{R}^i\Gamma(X,\mathscr{F})=0$ 对 $0<i\le k$。$k=0$ 时结果空洞地为真；于是设我们对一切 $0<k'<k$（把 $X$ 换成\emph{任何}一个\trAffineScheme）已知结果。设 $\alpha\in\mathrm{R}^k\Gamma(X,\mathscr{F})$。我们希望证明 $\alpha=0$。选取一个\emph{由 $\OO_X$-模构成的}内射分解
\[
\begin{tikzcd}[row sep=1.8em]
 & \vdots & \\
 & \mathscr{Q}_1 \ar[u, "d_1"] & \\
 & \mathscr{Q}_0 \ar[u, "d_0"] & \\
 & \mathscr{F} \ar[u] & \\
 & 0 \ar[u] &
\end{tikzcd}
\]
于是 $\alpha$ 有一个代表元 $\alpha'\in\mathscr{Q}_k(X)=\Gamma(X,\mathscr{Q}_k)$，使得 $d\alpha'=0$。由于内射分解是正合的，$\alpha'$ 局部地是一个边缘。换句话说，在 $X$ 中任意点 $\mathfrak{p}$ 的一个开邻域内，存在开集 $V_\mathfrak{p}$，使得对某个 $\alpha''\in\mathscr{Q}_{k-1}(V_\mathfrak{p})$ 有 $\alpha'|_{V_\mathfrak{p}}=d\alpha''$。（你确信这是真的吗？回想：取层态射的余核需要层化；见\mbox{命题~2.6.1}。）必要时缩小 $V_\mathfrak{p}$，我们不妨设 $V_\mathfrak{p}$ 是\trAffine{}的。由 $X$ 的拟紧性，我们可以选出有限多个 $V_\mathfrak{p}$ 覆盖 $X$。把这些重新命名为 $\mathfrak{U}_i$，于是我们有了 $X$ 的一个\trAffine{}覆盖。考虑 $X$ 相应于\emph{这个}\trAffine{}覆盖的切赫覆盖（\emph{不}是你可能以为我们会用的那个——即拿 $X$ 自身当作覆盖——而是一个为我们这个特定的 $\alpha$ 量身定制的仿射覆盖），把\trDoubleComplex{} (23.5.2.1) 作为某个谱序列的 $E_0$ 项。

先考虑右行定向。如 \S23.5.2 中的论证那样，谱序列在 $E_2$ 处收敛，此处我们处处得到 $0$，唯独导出函子上同调出现在第 $0$ 列。

接着，从头再来，选取上行定向。在 $E_1$ 页上，第 $0$ 行如 (23.5.2.2) 中那样是切赫复形。由归纳假设，$1$ 与 $k-1$ 之间的一切行都是 $0$，但对更高的行我们尚未所知。由于我们感兴趣的是第 $k$ 个导出函子，我们把注意力放在第 $k$ 条反对角线（$E_\bullet^{p,k-p}$）上。这条反对角线上仅有的可能非零的项是 $E_1^{k,0}$ 与 $E_1^{0,k}$。我们先看底行上的那一项 $E_1^{k,0}=\prod_{|I|=k}\Gamma(\mathfrak{U}_I,\mathscr{F})$，它是切赫复形的一部分：
\[
\cdots \longrightarrow \prod_{|I|=k-1}\Gamma(\mathfrak{U}_I,\mathscr{F}) \longrightarrow \prod_{|I|=k}\Gamma(\mathfrak{U}_I,\mathscr{F}) \longrightarrow \prod_{|I|=k+1}\Gamma(\mathfrak{U}_I,\mathscr{F}) \longrightarrow \cdots.
\]
但我们已经验证过：仿射概形上拟凝聚层的切赫上同调消失——这是我们所用「$\mathscr{F}$ 拟凝聚」的唯一一处。于是这一项由 $E_2$ 页起消失（即对 $i\ge 2$ 有 $E_i^{k,0}=0$）。

于是 $E_1$ 的第 $k$ 条反对角线上唯一有趣的项是 $E_1^{0,k}$，它是
\begin{equation*}\tag{23.5.4.1}\label{eq:agree-antidiag}
\prod_i\mathscr{Q}_{k-1}(\mathfrak{U}_i) \longrightarrow \prod_i\mathscr{Q}_k(\mathfrak{U}_i) \longrightarrow \prod_i\mathscr{Q}_{k+1}(\mathfrak{U}_i),
\end{equation*}
的\trHomology，即 $\prod_i\mathrm{R}^k\Gamma(\mathfrak{U}_i,\mathscr{F})$（用到\mbox{预备练习~23.5.A}：它说 $\mathscr{Q}_j|_{\mathfrak{U}_i}$ 在 $\mathfrak{U}_i$ 上内射，故它们可用来计算 $\mathrm{R}^k(\Gamma(\mathfrak{U}_i,\mathscr{F}))$）。于是 $E_2^{0,k}$ 是
\[
0 \longrightarrow \prod_i\mathrm{R}^k\Gamma(\mathfrak{U}_i,\mathscr{F}) \longrightarrow \prod_{i,j}\mathrm{R}^k\Gamma(\mathfrak{U}_{ij},\mathscr{F}),
\]
的\trHomology，此后进出 $E_\bullet^{0,k}$ 诸项的其余微分将都是 $0$，因为那些箭头的源与靶都是 $0$。现在考虑我们的 $\alpha\in\mathrm{R}^k\Gamma(X,\mathscr{F})$ 在 $\prod_i\mathrm{R}^k\Gamma(\mathfrak{U}_i,\mathscr{F})$ 中的提升 $\alpha'$。它在 (23.5.4.1) 的\trHomology{}中的像是\emph{零}——这正是我们起初选取覆盖 $\mathfrak{U}_i$ 的方式！于是 $\alpha=0$，正如所愿，证明完成。\hfill$\square$

\textbf{23.5.E.}\ \extrahard\ \textbf{（练习）}\quad 这个证明还没有完全收尾。我们有类 $\alpha\in\mathrm{R}^k\Gamma(X,\mathscr{F})$，而我们已经把 $\mathrm{R}^k\Gamma(X,\mathscr{F})$ 解读为
\[
\ker\Bigl(\prod_i\mathrm{R}^k\Gamma(\mathfrak{U}_i,\mathscr{F}) \longrightarrow \prod_{i,j}\mathrm{R}^k\Gamma(\mathfrak{U}_{ij},\mathscr{F})\Bigr).
\]
我们有两个映射 $\mathrm{R}^k\Gamma(X,\mathscr{F})\to\mathrm{R}^k\Gamma(\mathfrak{U}_i,\mathscr{F})$：一个来自自然的限制（在它之下我们可以看出 $\alpha$ 的像是零），另一个来自真正的谱序列机器。请验证：它们是同一个映射。（可能的提示：按所用的谱序列定向，$E_\infty^{0,k}$ 项确实是\trDoubleComplex{}的\trHomology{}的一个商，故该映射的方向是对的。）

\textbf{23.5.5.}\ \hard\ \textbf{收束线头。}\quad

\textbf{23.5.F.（重要练习）}\quad 请陈述并证明\mbox{定理~23.5.1} 到\trHigherPushforward{} $\mathrm{R}^i\pi_\ast$ 的推广，其中 $\pi\colon X\to Y$ 是概形的一个拟紧分离\trMorphism。

\textbf{23.5.G.（练习）}\quad 证明\mbox{定理~23.5.1} 的同构对 $\mathscr{F}$ 是\trFunctoriality{}的；也就是说，给定一个态射 $\mathscr{F}\to\mathscr{G}$，则方块
\[
\begin{tikzcd}[row sep=2.4em, column sep=3.4em]
\mathrm{H}^i(X,\mathscr{F}) \ar[d] & \mathrm{R}^i\Gamma(X,\mathscr{F}) \ar[l, <->] \ar[d] \\
\mathrm{H}^i(X,\mathscr{G}) & \mathrm{R}^i\Gamma(X,\mathscr{G}) \ar[l, <->]
\end{tikzcd}
\]
交换，其中诸水平箭头是\mbox{定理~23.5.1} 的同构，而诸竖直箭头来自 $\mathrm{H}^i$ 与 $\mathrm{R}^i\Gamma$ 的函子性。（提示：「谱序列在 $E_0$ 处函子」，这一点从构造中即可容易看出，尽管我们没有明说。参看 \S1.6.7，那里论及谱序列的函子性。）

\textbf{23.5.H.（练习）}\quad 证明：\mbox{定理~23.5.1} 的诸同构诱导出诸\trLongExactSequence{}的同构。

\textbf{23.5.6.}\ \emph{\trhRemark。}\quad 如果你愿意，你可以用上面的论证来证明\trLeray{}的下述定理。设 $\mathscr{F}$ 是拓扑空间 $X$ 上的一个由阿贝尔群构成的层，而 $\{\mathfrak{U}_i\}$ 是 $X$ 的一个覆盖，使得 $\mathscr{F}$ 在正次数处的（导出函子）上同调在诸 $\mathfrak{U}_i$ 的每个有限交上都消失。那么 $\mathscr{F}$ 的上同调可以用覆盖 $\{\mathfrak{U}_i\}$ 的切赫上同调来计算；无须像一般切赫上同调的情形那样，过渡到一切覆盖的余极限。

\textbf{23.5.7.}\ \emph{不重要的注记：在 $\mathit{QCoh}_X$ 中而非 $\mathit{Mod}_{\OO_X}$ 中工作。}\quad 在定义 $X$ 上拟凝聚层的导出函子时，我们本来也可以试着在拟凝聚层的范畴 $\mathit{QCoh}_X$ 自身中工作，而不是在更大的范畴 $\mathit{Mod}_{\OO_X}$ 中。出于若干理由，那样做将需要更多气力。不难证明：若 $X$ \trNoether，则 $\mathit{QCoh}_X$ 有足够多内射对象（比如见 [Ha1, Exer.~III.3.6(a)]）。由于在 $\mathit{QCoh}$ 中我们没有「\trExtensionByZero」（\S23.4.7），「开集 $\mathfrak{U}$ 上的内射拟凝聚层限制到更小的开子集 $V$ 上仍是内射拟凝聚层」（\mbox{练习~23.5.A} 的类似物）与「内射拟凝聚层是\trFlasque{}的」（\mbox{练习~23.4.H} 的类似物）这两个证明都会更难。你可以借此证明：在 $\mathit{QCoh}$ 中算出的 $\mathrm{H}^i$ 与 $\mathrm{R}^i\pi_\ast$，与在 $\mathit{Mod}_{\OO}$ 中算出者一致（一旦你把命题说精确）。内射的 $\mathit{QCoh}_X$ 对象（$X$ \trNoether）在 $\mathit{Mod}_{\OO_X}$ 中也是内射的，这是真的，但其证明需要功夫（见 [Mur, Prop.~6.8]）。

对\emph{任何}概形 $X$，$\mathit{QCoh}_X$ 都有足够多内射对象，这也是真的，但那要难得多；见 [EE]。而从上一段可以清楚看到，「有足够多内射对象」只是我们想要的东西的开端。

\textbf{23.5.8.}\ \emph{不重要的注记。}\quad \mbox{定理~23.5.1} 蕴含：若 $\pi\colon X\to Y$ 拟紧且分离，则 $\mathrm{R}^i\pi_\ast$ 把 $\mathit{QCoh}_X$ 送到 $\mathit{QCoh}_Y$（通过展示与切赫上同调的一个同构）。若 $X$ 与 $Y$ \trNoether，则借助\mbox{不重要注记~23.5.7} 的想法，「分离」这一假设可以放宽为「拟分离」。但由于从\trNoether{}概形出发的态射自动是拟分离的（\mbox{容易练习~8.3.B}(b)），我们可以把这个多余的「拟分离」假设去掉。这听起来令人兴奋，但它远没有分离的情形那么有用，因为没有切赫上同调，就很难算出任何东西。



